%% PRACA DYPLOMOWA. Podstawa: szablon WIT PWr z 25.02.2026.
%% Kompilacja: pdflatex, biber, pdflatex, pdflatex (szczegóły w README.md).
%% Wszystkie pliki w UTF-8.

\documentclass[12pt, a4paper, twoside, titlepage]{report}

%%% Copyright (C) [2026] [PWr * WIT * Wydziałowa Komisja d/s Jakości Kształcenia]
%%% Zmodyfikowany i zoptymalizowany (2026)
%%% ---------------------------------------------------------------------------

%%% KONFIGURACJA EKSTRAKCJI TEKSTU (PDF COPY-PASTE)
\input{glyphtounicode} 
\pdfgentounicode=1     % Wymagane przez JSA (Jednolity System Antyplagiatowy)

%%% 1. KODOWANIE I JĘZYK
\RequirePackage[T1]{fontenc}
%\RequirePackage[utf8]{inputenc} % Opcjonalne w nowych systemach (domyślne)
\RequirePackage[english,polish]{babel} % Kolejność eng -> pl
\RequirePackage{csquotes} % Wymagane przez biblatex

%%% 2. MATEMATYKA I CZCIONKI
\RequirePackage{amsmath}
\RequirePackage{mathtools}
\RequirePackage{amsthm}

\RequirePackage{newtxtext} 
\RequirePackage{newtxmath}
\RequirePackage{microtype} % Poprawia optyczny wygląd tekstu

%%% 3. GEOMETRIA I GRAFIKA
\RequirePackage{geometry}
\geometry{
  a4paper,
  top=25mm, bottom=25mm,
  inner=25mm, outer=25mm,
  bindingoffset=5mm
}
\RequirePackage{graphicx}
\RequirePackage{xcolor}
\RequirePackage{float} % Dla opcji pozycjonowania [H]

%%% 4. NARZĘDZIA POMOCNICZE
\RequirePackage{etoolbox}

%%% 5. FORMATOWANIE TEKSTU I NAGŁÓWKÓW
\RequirePackage{titlesec}
\RequirePackage{caption}
\RequirePackage{subcaption}
\RequirePackage{setspace}

\setlength{\parindent}{0.7cm}
\setstretch{1.15} 

% Formatowanie nagłówków (Rozdział, Podrozdział)
\titleformat{\chapter}[display]
  {\bfseries\Large}
  {\chaptertitlename\ \thechapter}{0.5ex}   
  {\huge}
\titlespacing*{\chapter}{0pt}{-20pt}{20pt}

\titleformat{\section}{\bfseries\Large}{\thesection.}{1em}{}
\titlespacing*{\section}{0pt}{2.5ex plus 1ex minus .2ex}{1.3ex plus .2ex}

\titleformat{\subsection}{\bfseries\large}{\thesubsection.}{1em}{}
\titlespacing*{\subsection}{0pt}{2.25ex plus 1ex minus .2ex}{1.0ex plus .2ex}

% Numeracja ciągła wewnątrz rozdziałów
\counterwithin{figure}{chapter}
\counterwithin{table}{chapter}
\numberwithin{equation}{chapter}

% Podpisy (Captions) - pogrubione etykiety
\captionsetup{font=small, labelfont=bf, labelsep=period}

%%% 6. ALGORYTMY I LISTINGI KODU
\RequirePackage{algorithm}
\RequirePackage{algpseudocode}
\algrenewcommand\algorithmiccomment[1]{\hfill\(\triangleright\) #1}
\floatname{algorithm}{Algorytm} 

\RequirePackage{listings}
\definecolor{codebg}{RGB}{250,250,250}
\definecolor{codekw}{RGB}{0,64,160}
\definecolor{codestr}{RGB}{0,128,0}
\definecolor{codecmt}{RGB}{128,128,128}

\lstset{
  basicstyle=\ttfamily\small,
  backgroundcolor=\color{codebg},
  keywordstyle=\bfseries\color{codekw},
  stringstyle=\color{codestr},
  commentstyle=\itshape\color{codecmt},
  numbers=left,
  numberstyle=\tiny\color{gray},
  stepnumber=1,
  numbersep=8pt,
  tabsize=2,
  showstringspaces=false,
  breaklines=true,
  frame=lines,
  captionpos=b,
  % Mapowanie polskich znaków dla listings (niezbędne przy tym pakiecie)
  literate={ą}{{\k{a}}}1 {ć}{{\'c}}1 {ę}{{\k{e}}}1 {ł}{{\l{}}}1 {ń}{{\'n}}1 {ó}{{\'o}}1 {ś}{{\'s}}1 {ż}{{\.z}}1 {ź}{{\'z}}1
           {Ą}{{\k{A}}}1 {Ć}{{\'C}}1 {Ę}{{\k{E}}}1 {Ł}{{\L{}}}1 {Ń}{{\'N}}1 {Ó}{{\'O}}1 {Ś}{{\'S}}1 {Ż}{{\.Z}}1 {Ź}{{\'Z}}1
}
\renewcommand{\lstlistingname}{Listing}
\renewcommand{\lstlistlistingname}{Spis listingów}

% Styl Haskell 
\definecolor{darkgreen}{rgb}{0,0.5,0}
\lstdefinestyle{haskellStyle}{
    language=Haskell,
    commentstyle=\color{darkgreen},
    keywordstyle=\color{blue}\bfseries,
    basicstyle=\ttfamily\small,
    breaklines=true,
    stringstyle=\color{red}
}

%%% 7. HYPERREF I CLEVEREF (Musi być na końcu)
\RequirePackage[
    pdfusetitle,
    colorlinks=true,
    linkcolor=black,
    citecolor=blue,
    urlcolor=blue,
    bookmarksnumbered=true
]{hyperref}

\RequirePackage[nameinlink]{cleveref}

% Tłumaczenia dla cleveref
\crefname{figure}{Rys.}{Rys.}
\Crefname{figure}{Rysunek}{Rysunki}
\crefname{table}{Tab.}{Tab.}
\Crefname{table}{Tabela}{Tabele}
\crefname{equation}{równ.}{równ.}
\Crefname{equation}{Równanie}{Równania}
\crefname{listing}{listing}{Listingi}
\crefname{algorithm}{algorytm}{algorytmy}
\Crefname{algorithm}{Algorytm}{Algorytmy}

%%%%%%%%%%%
%% MAKRA %%
%%%%%%%%%%%

% Definicja rodzaju pracy 
\newcommand{\stRodzajPracy}{%
  \ifnum\stCzyMgr=1 MAGISTERSKA\else INŻYNIERSKA\fi
}

\newcommand{\StronaTytulowa}{
  \pagenumbering{roman}
  \sffamily
  \begin{titlepage}  
    \newgeometry{top=2.5cm, bottom=2.5cm, left=2.5cm, right=2.5cm}
    
    \begin{center}
      {\huge\bfseries Politechnika Wrocławska \par}
      \vspace{0.4cm}
      {\Large Wydział Informatyki i Telekomunikacji \par}
    \end{center}
    \par\noindent\rule{\textwidth}{1pt}\par 
    \vspace{0.5cm} 
    
    {\large
    \begin{tabular}{@{}ll}
      Kierunek: & \textbf{\stKierunek\ (\stKodKierunku)} \\
      \ifdefempty{\stSpecjalnosc}{}{Specjalność: & \textbf{\stSpecjalnosc\ (\stKodSpecjalnosci)}}
    \end{tabular}
    }
    
    \vspace*{\fill}
    
    \begin{center}
      {\huge\bfseries PRACA DYPLOMOWA \par}
      \vspace{0.3cm}    
      {\huge \stRodzajPracy \par}
      \vspace{2.0cm}
      {\Large \bfseries \stTytulPL \par} 
      \vspace{1cm}
      {\Large \stAutor \par}
      \vspace{2cm}
      {\large Opiekun pracy:\par \textbf{\stOpiekun} \par}
    \end{center}
    
    \vspace*{\fill}
    
    {\raggedright
      \small Słowa kluczowe: \textit{\stSlowaKluczowe}\par
      \vspace{0.5em}
      \par\noindent\rule{\textwidth}{1pt}\par 
    }
    
    \begin{center}
      \Large Wrocław \stRok
    \end{center}
  
    \restoregeometry
  \end{titlepage}
  \normalfont
  \pagestyle{empty}
}

\newenvironment{abstracts}{}{\selectlanguage{polish}\clearpage}

\newcommand{\AbstractPL}{
  \begin{center}{\Large\bfseries Streszczenie \par}\end{center}\vspace{0.5em}
}

\newcommand{\AbstractENG}{
  \selectlanguage{english}
  \begin{center}{\Large\bfseries Abstract \par}\end{center}\vspace{0.5em}
}

\newcommand{\SpisTresci}{
  \clearpage
  \addtocontents{toc}{\protect\thispagestyle{empty}} 
  \pagestyle{empty} 
  \tableofcontents
  \clearpage
  \pagestyle{empty}
}

%% WSTAWIENIE METADANYCH DO PLIKU PDF
\newcommand{\SetMetaData}{%
  \AtBeginDocument{%  
    \hypersetup{%
      pdfinfo = {%
        Title    = {\stTytulPL},
        Author   = {\stAutor},
        Subject  = {Praca dyplomowa \stRodzajPracy},  
        Keywords = {\stSlowaKluczowe}, 
        ModDate  = {\pdfcreationdate}, 
        Creator  = {pdftex}
      }
    }
  }
}
%%% KONIEC PLIKU SETTINGS.TEX         % plik Wydziału, bez zmian
%%% DODATKI DO SZABLONU WIT
%%% Plik settings.tex zostaje bez zmian; wszystko, co dodane, jest tutaj.

%%% 1. TABELE
\RequirePackage{array}
\RequirePackage{booktabs}
\RequirePackage{tabularx}
\RequirePackage{multirow}
\RequirePackage{longtable}
\newcolumntype{L}{>{\raggedright\arraybackslash}X}
\newcolumntype{C}{>{\centering\arraybackslash}X}

% Zalecenia edytorskie Wydziału: tytuł tabeli centralnie nad tabelą, podpis rysunku
% centralnie pod rysunkiem, tabela 10 pkt. Przy klasie 12 pkt to \footnotesize: każda tabela
% ma je wpisane sama, bo LaTeX zeruje rozmiar czcionki na początku wstawki.
\captionsetup{justification=centering}
\captionsetup[table]{position=top, skip=6pt}
\AtBeginEnvironment{longtable}{\footnotesize}

% Domyślne położenie: klocki nie podają go same, więc w katalogu mogą stać w miejscu.
\floatplacement{table}{htbp}
\floatplacement{figure}{htbp}

%%% 2. POLSKIE NAZWY DLA \cref (settings.tex definiuje tylko rysunki, tabele, równania, listingi, algorytmy)
\crefname{chapter}{rozdz.}{rozdz.}
\Crefname{chapter}{Rozdział}{Rozdziały}
\crefname{section}{podrozdz.}{podrozdz.}
\Crefname{section}{Podrozdział}{Podrozdziały}
\crefname{subsection}{podrozdz.}{podrozdz.}
\Crefname{subsection}{Podrozdział}{Podrozdziały}
\crefname{appendix}{zał.}{zał.}
\Crefname{appendix}{Załącznik}{Załączniki}
\crefname{subfigure}{Rys.}{Rys.}
\Crefname{subfigure}{Rysunek}{Rysunki}
\Crefname{listing}{Listing}{Listingi}
\crefname{lstlisting}{listing}{listingi}
\Crefname{lstlisting}{Listing}{Listingi}

%%% 3. STYL KODU C#
%%% literate zastępuje ustawienie z settings.tex, więc są tu też polskie znaki.
\lstdefinestyle{csharp}{
  language=[Sharp]C,
  morekeywords={var, get, set, yield, where, nameof, unsafe, stackalloc, ref},
  % 10 pkt jak w tabelach; długi wiersz łamie się na spacji, a ciąg dalszy ma strzałkę.
  basicstyle=\ttfamily\footnotesize,
  breakatwhitespace=true,
  postbreak=\mbox{\textcolor{gray}{$\hookrightarrow$}\space},
  literate={ą}{{\k{a}}}1 {ć}{{\'c}}1 {ę}{{\k{e}}}1 {ł}{{\l{}}}1 {ń}{{\'n}}1 {ó}{{\'o}}1 {ś}{{\'s}}1 {ż}{{\.z}}1 {ź}{{\'z}}1
           {Ą}{{\k{A}}}1 {Ć}{{\'C}}1 {Ę}{{\k{E}}}1 {Ł}{{\L{}}}1 {Ń}{{\'N}}1 {Ó}{{\'O}}1 {Ś}{{\'S}}1 {Ż}{{\.Z}}1 {Ź}{{\'Z}}1
           {×}{{$\times$}}1 {–}{{\textendash}}1 {−}{{$-$}}1 {·}{{$\cdot$}}1 {°}{{$^\circ$}}1 {…}{{$\ldots$}}1
           {→}{{$\rightarrow$}}1 {↔}{{$\leftrightarrow$}}1 {≈}{{$\approx$}}1 {κ}{{$\kappa$}}1 {Σ}{{$\Sigma$}}1
           {„}{{\quotedblbase}}1 {”}{{\textquotedblright}}1 {•}{{$\bullet$}}1 {■}{{$\blacksquare$}}1 {✓}{{$\checkmark$}}1
}

%%% 4. PRZEŁĄCZNIK SZKICU I MIEJSCA NA TREŚĆ
%%% \szkictrue albo \szkicfalse ustawia się w main.tex.
\newif\ifszkic
\definecolor{szkicramka}{RGB}{176,112,0}
\definecolor{szkictlo}{RGB}{255,246,224}
\newcounter{szkicwskazowki}
\newcounter{szkicpola}
\newcounter{szkicobrazy}
\newcounter{szkiclistingi}

\newcommand{\szkicRamka}[1]{%
  \par\smallskip\noindent
  \fcolorbox{szkicramka}{szkictlo}{%
    \parbox{\dimexpr\linewidth-2\fboxsep-2\fboxrule\relax}{%
      \raggedright\footnotesize\sffamily #1}}%
  \par\smallskip}

% Wskazówka w ramce; po \szkicfalse znika, a \cite w niej przestaje się liczyć do bibliografii.
\newcommand{\CoTuMaByc}[1]{%
  \stepcounter{szkicwskazowki}%
  \ifszkic\szkicRamka{\textbf{Co tu ma być:} #1}\fi}

% Pole do wpisania w tekście, tabeli albo pseudokodzie; po \szkicfalse jest puste.
\newcommand{\Wpisz}[1]{%
  \stepcounter{szkicpola}%
  \ifszkic\textcolor{szkicramka}{[#1]}\fi}

% Pusta komórka tabeli na wynik.
\newcommand{\wyn}{\Wpisz{?}}

% Nazwa kolumny CSV albo pliku z podkreśleniami.
\newcommand{\kol}[1]{\texttt{\detokenize{#1}}}

% Rozmiar siatki: E × N × piętra.
\newcommand{\siatka}[3]{#1\,$\times$\,#2\,$\times$\,#3}

% Obraz z folderu figures, bez rozszerzenia. Gdy pliku nie ma, stoi ramka zastępcza.
% [szerokość]{plik}{co ma przedstawiać}
\newcommand{\MiejsceNaObraz}[3][0.85\linewidth]{%
  \IfFileExists{figures/#2.pdf}{\includegraphics[width=#1]{figures/#2.pdf}}{%
  \IfFileExists{figures/#2.png}{\includegraphics[width=#1]{figures/#2.png}}{%
  \IfFileExists{figures/#2.jpg}{\includegraphics[width=#1]{figures/#2.jpg}}{%
    \stepcounter{szkicobrazy}%
    \fbox{\parbox[c][\dimexpr#1*9/16\relax][c]{\dimexpr#1-2\fboxsep-2\fboxrule\relax}{%
      \centering
      \ifszkic
        \footnotesize\sffamily\textbf{figures/\detokenize{#2}}\par\smallskip #3%
      \fi}}%
  }}}}

% Cały rysunek: [szerokość]{plik}{podpis}{etykieta}{co ma przedstawiać}
\newcommand{\MiejsceNaRysunek}[5][0.85\linewidth]{%
  \begin{figure}
    \centering
    \MiejsceNaObraz[#1]{#2}{#5}
    \caption{#3}
    \label{#4}
  \end{figure}}

% Listing z folderu listingi, bez rozszerzenia .cs: {plik}{podpis}{etykieta}{skąd wziąć kod}
\newcommand{\MiejsceNaListing}[4]{%
  \IfFileExists{listingi/#1.cs}{%
    \lstinputlisting[style=csharp, caption={#2}, label={#3}]{listingi/#1.cs}%
  }{%
    \stepcounter{szkiclistingi}%
    \ifszkic\szkicRamka{\textbf{Listing:} zapisz fragment jako \kol{listingi/#1.cs}. #4}\fi
    \lstinputlisting[style=csharp, caption={#2}, label={#3}]{listingi/brak.cs}%
  }}

% Po kompilacji: ile miejsc zostało do uzupełnienia (ostrzeżenie w dzienniku kompilacji).
\AtEndDocument{%
  \ifnum\numexpr\value{szkicwskazowki}+\value{szkicpola}+\value{szkicobrazy}+\value{szkiclistingi}\relax>0
    \PackageWarningNoLine{szablon pracy}{Do uzupelnienia: wskazowki \arabic{szkicwskazowki}, pola \arabic{szkicpola}, rysunki \arabic{szkicobrazy}, listingi \arabic{szkiclistingi}}%
  \fi}

%%% 5. ZALECENIA EDYTORSKIE WYDZIAŁU
%%% Szablon WIT ma większe nagłówki i interlinię 1,15. Zalecenia edytorskie: tytuł rozdziału 14 pkt,
%%% podrozdziału 13 pkt, niższego stopnia 12 pkt, interlinia pojedyncza. Włączenie: jedna linia w main.tex.
\newcommand{\WlaczZaleceniaEdytorskie}{%
  \setstretch{1}%
  \titleformat{\chapter}[display]{\bfseries\fontsize{14}{17}\selectfont}{\chaptertitlename\ \thechapter}{0.5ex}{}%
  \titleformat{\section}{\bfseries\fontsize{13}{16}\selectfont}{\thesection.}{1em}{}%
  \titleformat{\subsection}{\bfseries\fontsize{12}{15}\selectfont}{\thesubsection.}{1em}{}%
}
          % tabele, styl C#, makra miejsc na treść
%%% STRONA TYTUŁOWA WEDŁUG WZORU WYDZIAŁU pd_inz_pl.doc
%%% Marginesy wzoru: górny 25 mm, dolny 17,5 mm, lewy 35 mm, prawy 25 mm.
%%% Odległości w punktach bp odczytano z pliku PDF wyeksportowanego ze wzoru.
%%% Wzór używa czcionki Garamond; tu EB Garamond (pakiet ebgaramond, jest w TeX Live i na Overleaf).
%%% Gdy pakietu nie ma, strona powstaje czcionką główną pracy.
%%% Długi tytuł zabiera tylko pusty odstęp nad słowami kluczowymi; dół strony stoi w miejscu.

\newcommand{\stCzcionka}{%
  \IfFileExists{T1EBGaramond-LF.fd}%
    {\fontfamily{EBGaramond-LF}\selectfont}%
    {\rmfamily}}

% Wiersz: rozmiar czcionki w pkt, odległość linii bazowej od poprzedniego wiersza w bp, treść.
\newcommand{\stWiersz}[3]{{\fontsize{#1}{#2bp}\selectfont #3\par}}

% Kreska 0,48 bp, nieco szersza od tekstu jak we wzorze; #1 = odległość od linii bazowej wiersza nad nią.
\newcommand{\stKreska}[1]{%
  \par
  \vskip\dimexpr#1-\prevdepth\relax
  \nointerlineskip
  \moveleft1.43bp\vbox{\hrule height0.48bp width\dimexpr\linewidth+2.86bp\relax}%
}

% Pierwsza litera 26 pkt, reszta wiersza 24 pkt.
\newcommand{\stInicjal}[2]{{\fontsize{26}{29.3bp}\selectfont #1}#2}

% Kierunek i specjalność: etykieta, położenie wartości od lewego marginesu, wartość.
\newlength{\stSzerokosc}
\newcommand{\stPole}[3]{%
  \settowidth{\stSzerokosc}{\fontsize{14}{15.7bp}\selectfont #1\ }%
  \ifdim\stSzerokosc<#2\relax\setlength{\stSzerokosc}{#2}\fi
  \stWiersz{14}{15.7}{\raggedright\hangindent=\stSzerokosc\hangafter=1\relax
    \makebox[\stSzerokosc][l]{#1}\textbf{#3}}}

\renewcommand{\StronaTytulowa}{%
  \pagenumbering{roman}%
  \begin{titlepage}%
    \newgeometry{top=25mm, bottom=17.5mm, left=35mm, right=25mm}%
    \setstretch{1}%
    \stCzcionka
    \setlength{\parindent}{0pt}%
    \setlength{\parskip}{0pt}%
    \setlength{\topskip}{17.25bp}%
    \lineskiplimit=-1000pt\relax
    {\centering
      \stWiersz{20}{22.5}{\bfseries Politechnika Wrocławska}%
      \stWiersz{16}{19.1}{\bfseries Wydział Informatyki i Telekomunikacji}%
    \par}%
    \stKreska{5.02bp}%
    \vskip19bp\relax
    \stPole{Kierunek:}{106.25bp}{\stKierunek\ (\stKodKierunku)}%
    \ifdefempty{\stSpecjalnosc}%
      {\vskip15.7bp\relax}%
      {\stPole{Specjalność:}{70.85bp}{\stSpecjalnosc\ (\stKodSpecjalnosci)}}%
    {\centering
      \vskip59.9bp\relax
      \stWiersz{24}{29.3}{\hspace*{7.1bp}\stInicjal{P}{RACA} \stInicjal{D}{YPLOMOWA}}%
      \stWiersz{24}{29.3}{\hspace*{7.1bp}\ifnum\stCzyMgr=1 \stInicjal{M}{AGISTERSKA}\else\stInicjal{I}{NŻYNIERSKA}\fi}%
      \vskip74.15bp\relax
      \stWiersz{18}{20.25}{\bfseries\stTytulPL}%
      \vskip36.5bp\relax
      \stWiersz{16}{18}{\stAutor}%
      \vskip56.65bp\relax
      \stWiersz{14}{15.75}{Opiekun pracy}%
      \stWiersz{14}{15.8}{\bfseries\stOpiekun}%
    \par}%
    \vfill
    \stWiersz{12}{13.5}{\raggedright Słowa kluczowe: \stSlowaKluczowe}%
    \stKreska{39.98bp}%
    {\centering\stWiersz{14}{12.14}{WROCŁAW \stRok}\par}%
    \vskip\dimexpr47.09bp-\prevdepth\relax
    \restoregeometry
  \end{titlepage}%
  \normalfont
  \pagestyle{empty}%
}
  % strona tytułowa według pd_inz_pl.doc;
                         % zakomentuj tę linię, aby wrócić do strony tytułowej z settings.tex

%% PRZEŁĄCZNIK SZKICU
%% \szkictrue:  widać ramki ze wskazówkami, puste komórki, rysunki zastępcze i katalog klocków
%% \szkicfalse: wskazówki i katalog znikają, zostaje wygląd końcowy
\szkictrue

%% Zalecenia edytorskie Wydziału (nagłówki 14/13/12 pkt, interlinia pojedyncza).
%% Bez tej linii obowiązuje wygląd z szablonu WIT.
%\WlaczZaleceniaEdytorskie

%% DANE PRACY
\newcommand{\stKierunek}{ Informatyka Techniczna}
\newcommand{\stKodKierunku}{ITE} %% kody: skroty_kier_pd.xls na stronie Wydziału
\newcommand{\stSpecjalnosc}{Grafika i Multimedia} %% pusty wpis = brak specjalności, wiersz znika
\newcommand{\stKodSpecjalnosci}{IGM} %% kody: skroty_kier_pd.xls na stronie Wydziału
\newcommand{\stCzyMgr}{0} % 1 = Magisterska, 0 = Inżynierska
%% Tytuł musi być identyczny z tytułem w Archiwum Prac Dyplomowych.
\newcommand{\stTytulPL}{Analiza porównawcza i optymalizacja algorytmu Wave Function Collapse w procesie proceduralnej generacji środowisk 3D w środowisku Unity}
%% Tytuł angielski: tłumaczenie robocze; musi być identyczny z tytułem angielskim w Archiwum Prac Dyplomowych.
\newcommand{\stTytulAng}{Comparative analysis and optimization of the Wave Function Collapse algorithm in procedural generation of 3D environments in the Unity engine}
\newcommand{\stAutor}{Jakub Żłobiński}
\newcommand{\stOpiekun}{dr inż. Dominik Żelazny}
\newcommand{\stSlowaKluczowe}{Proceduralna generacja, Unity, Wave Function Collapse}
\newcommand{\stRok}{2026}  % rok złożenia pracy

\SetMetaData % metadane pliku PDF

%% WŁASNE SKRÓTY
\newcommand{\NN}{\mathbb{N}} % liczby naturalne
\newcommand{\RR}{\mathbb{R}} % liczby rzeczywiste

\newtheorem{theorem}{Twierdzenie}[chapter]
\newtheorem{lemma}{Lemat}[chapter]
\newtheorem{corollary}{Wniosek}[chapter]
\newtheorem{definition}{Definicja}[chapter]

%% BIBLIOGRAFIA: biblatex i program biber
\RequirePackage[
    backend=biber,
    style=numeric,
    sorting=none,     % kolejność występowania w tekście
    urldate=iso,
    seconds=true
]{biblatex}
\addbibresource{bibliography.bib}
% Długie adresy internetowe mogą się łamać także między literami i cyframi.
\setcounter{biburlnumpenalty}{100}
\setcounter{biburlucpenalty}{100}
\setcounter{biburllcpenalty}{100}

%% AKRONIMY
%%% LISTA AKRONIMÓW
%%% W pracy pisz \ac{wfc}. Pierwsze użycie daje pełną nazwę ze skrótem, następne sam skrót.
%%% Na liście pojawiają się tylko akronimy użyte w tekście.

\RequirePackage{acro}

\DeclareAcronym{ac3}{
  short = AC-3 ,
  long  = Arc Consistency Algorithm 3
}
\DeclareAcronym{ac4}{
  short = AC-4 ,
  long  = Arc Consistency Algorithm 4
}
\DeclareAcronym{ai}{
  short = AI ,
  long  = Artificial Intelligence
}
\DeclareAcronym{api}{
  short = API ,
  long  = Application Programming Interface
}
\DeclareAcronym{cpu}{
  short = CPU ,
  long  = Central Processing Unit
}
\DeclareAcronym{csp}{
  short = CSP ,
  long  = Constraint Satisfaction Problem
}
\DeclareAcronym{csv}{
  short = CSV ,
  long  = Comma-Separated Values
}
\DeclareAcronym{fbx}{
  short = FBX ,
  long  = Filmbox
}
\DeclareAcronym{gc}{
  short = GC ,
  long  = Garbage Collector
}
\DeclareAcronym{gpu}{
  short = GPU ,
  long  = Graphics Processing Unit
}
\DeclareAcronym{il2cpp}{
  short = IL2CPP ,
  long  = Intermediate Language To C++
}
\DeclareAcronym{json}{
  short = JSON ,
  long  = JavaScript Object Notation
}
\DeclareAcronym{mrv}{
  short = MRV ,
  long  = Minimum Remaining Values
}
\DeclareAcronym{pcg}{
  short = PCG ,
  long  = Procedural Content Generation
}
\DeclareAcronym{srp}{
  short = SRP ,
  long  = Scriptable Render Pipeline
}
\DeclareAcronym{urp}{
  short = URP ,
  long  = Universal Render Pipeline
}
\DeclareAcronym{wfc}{
  short = WFC ,
  long  = Wave Function Collapse
}

%%%%%%%%%%%%%%%%%%%%%%%%%%%%%%%%%%%%%%%%%%%
%% NIE RUSZAJ TEGO CO JEST NIŻEJ JEŚLI  %%%
%% NIE ZNASZ ZBYT DOBRZE LATEX'A        %%%
%%%%%%%%%%%%%%%%%%%%%%%%%%%%%%%%%%%%%%%%%%%
\newcommand{\ShowAcronyms}[1]{
  \begin{center}
    {\Large\bfseries Lista skrótów \par}
  \end{center}
  \vspace{0.5em}

  \ifboolexpr{
      test {\ifstrequal{#1}{lof}} or
      test {\ifstrequal{#1}{toc}} or
      test {\ifstrequal{#1}{description}}
    }{%
      \printacronyms[sort=true, heading=none, template=#1]
    }{%
      \printacronyms[sort=true, heading=none, template=lof]
     }%
  \clearpage
}


\begin{document}

%% NIE RUSZAJ NASTĘPNYCH TRZECH LINIJEK
\StronaTytulowa
%% ZACHOWAJ STRUKTURĘ TEGO PLIKU
%% TZN. abstrakty muszą być wewnątrz środowiska abstract
%% Polecenia \AbstractPL i \AbstractENG są zdefiniowane 
%% w pliku main,tex

\begin{abstracts}
  \AbstractPL
	Tutaj umieść zwięzłe streszczenie swojej pracy dyplomowej. 
	Powinno ono zawierać opis celu pracy, zastosowane metody oraz uzyskane wyniki.  \par
	
	Oba streszczenia, po polsku i po angielsku, powinny zmieścić się na jednej stronie.
		
	\AbstractENG
	Place a concise summary of your thesis here. 
	It should include a description of the purpose of the work, the methods used and the results obtained.
\end{abstracts}


\SpisTresci

\ShowAcronyms{lof}

\pagenumbering{arabic}
\pagestyle{plain} % dodane: w szablonie WIT numer strony był tylko na pierwszej stronie rozdziału

%% ROZDZIAŁY
% !TEX root = ../main.tex
%% Układ tego rozdziału pochodzi z szablonu WIT, a wskazówki z wymagań Wydziału.

\chapter{Wstęp}
\label{chap:wstep}

\CoTuMaByc{Wymagania Wydziału: cele pracy, motywacje jej podjęcia, krótki opis zawartości.}

\section{Cel i struktura pracy}
\label{sec:cel-pracy}

\CoTuMaByc{Cel pracy, zwięźle i precyzyjnie. Cel i zakres zgodne z opisem tematu pracy.}

\subsection*{Struktura pracy}

\CoTuMaByc{Po kolei, krótko, wszystkie rozdziały.}

\subsection*{Zastosowane technologie}

\CoTuMaByc{Języki, kompilatory, środowisko uruchomienia i testów, najważniejsze biblioteki. Tabela wersji: \cref{tab:technologie}.}

\section{Wykorzystanie AI}
\label{sec:wykorzystanie-ai}

\CoTuMaByc{Szablon WIT: tej części nie wolno pominąć, można ją tylko przenieść w inne miejsce pracy. Jakie narzędzia, do czego, których miejsc pracy to dotyczy. Jeśli AI generowało kod, trzeba o tym napisać tutaj i zaznaczyć te fragmenty w komentarzach kodu. Wytyczne Wydziału: \cite{WITWytyczneSI}. Zakres użycia uzgodnij z opiekunem. Wiersze tabeli wpisane według historii sesji Claude Code w tym projekcie (20.10--12.11.2026, jedna sesja, jeden model); do sprawdzenia i uzupełnienia o inne narzędzia, jeśli były używane, np. przy planie pracy, prototypie w Pythonie albo prezentacji.}

\begin{table}
  \centering
  \footnotesize
  \caption{Użyte narzędzia sztucznej inteligencji}
  \label{tab:narzedzia-ai}
  \begin{tabularx}{\linewidth}{@{}>{\raggedright\arraybackslash}p{3.1cm}LL@{}}
    \toprule
    Narzędzie i wersja & Do czego użyte & Których miejsc pracy dotyczy \\
    \midrule
    Claude Code 2.1.299 (Anthropic), model Claude Opus 5.5 & kod generatora, testów i aplikacji do badań (\kol{Assets/Scripts}); skrypty \kol{Analiza/analiza_wynikow.py} i \kol{Analiza/wykresy.py}; model \kol{Dach_Szczyt_Choragiew} w Blenderze i jego wpis w pliku reguł & listingi z kodu projektu; wykresy z wyników badań; \Wpisz{rozdziały o implementacji i badaniach} \\
    \addlinespace
    Claude Code 2.1.299 (Anthropic), model Claude Opus 5.5 & szablon LaTeX i strona tytułowa; tabele z danymi z kodu i ustawień projektu; zapis wzorów w LaTeX-u; wyszukanie i sformatowanie pozycji literatury; tłumaczenie tytułu na angielski & układ dokumentu; tabele z danymi projektu; wzory; załączniki; spis literatury; tytuł angielski \\
    \bottomrule
  \end{tabularx}
\end{table}

%% Tu dopisz swoje rozdziały, np.:
%% \include{chapters/rozdzial02}
%% \include{chapters/rozdzial03}

%% Magazyn gotowych tabel, rysunków i wzorów; widoczny tylko w szkicu.
\ifszkic
  %% KATALOG KLOCKÓW: nie jest rozdziałem pracy, widać go tylko w szkicu.
%% Klocek przenosisz przez wytnij i wklej z pliku w folderze klocki do swojego rozdziału.
%% Wtedy etykiety się nie dublują, a katalog pokazuje, co jeszcze zostało.

\chapter{Katalog klocków (do usunięcia)}
\label{chap:katalog}

\CoTuMaByc{Magazyn gotowych tabel, rysunków, wzorów, pseudokodów i listingów. Wytnij potrzebny blok z pliku w folderze \kol{klocki} i wklej do swojego rozdziału. Wymaganie Wydziału: każdą tabelę, każdy rysunek i każde źródło trzeba przywołać w tekście, np. przez \kol{\cref{tab:e1}}. Tabele powstałe z pomocą AI opisz w sekcji o AI (\cref{sec:wykorzystanie-ai}). Numery i odwołania przeliczą się same. Skróty w tekście wpisuj przez makro: \kol{\ac{wfc}} daje \ac{wfc}, a lista skrótów uzupełnia się sama.}

\begingroup
\floatplacement{table}{H}
\floatplacement{figure}{H}
\floatplacement{algorithm}{H}

\section{Fakty o projekcie}
%%% KLOCKI: FAKTY O PROJEKCIE
%%% Dane z kodu i zasobów projektu (stan: commit 4e10be9). Puste pola są do uzupełnienia.
%%% Przenoszenie: wytnij blok od \begin{table} do \end{table} razem z ramką nad nim.

\CoTuMaByc{Wersje pakietów z \kol{Packages/manifest.json} i \kol{ProjectSettings/ProjectVersion.txt}, wersja języka C\# z plików \kol{*.csproj}, Blender z nagłówka plików FBX, edytor kodu z ustawień Unity. Wersja kompilatora C++ będzie w dzienniku budowania (\kol{Editor.log}) po buildzie IL2CPP.}

\begin{table}
  \centering
  \footnotesize
  \caption{Użyte technologie i ich wersje}
  \label{tab:technologie}
  \begin{tabular}{@{}lll@{}}
    \toprule
    Składnik & Wersja & Zastosowanie \\
    \midrule
    Unity & 6000.6.2f1 & silnik i edytor \\
    Universal Render Pipeline (URP) & 17.6.0 & potok renderowania \\
    Burst & 2.0.0 & kompilacja zadań rozwiązywacza i scalania \\
    Collections & 6.6.0 & kontenery pamięci natywnej \\
    Mathematics & 1.4.0 & typy wektorowe \\
    Input System & 1.20.0 & sterowanie kamerą \\
    Test Framework & 1.8.0 & testy automatyczne \\
    C\# & 9.0 & język generatora \\
    IL2CPP (Unity) & 6000.6.2f1 & zamiana kodu C\# na C++ w aplikacji \\
    Microsoft C/C++ (MSVC) & 19.51 & kompilacja kodu C++ z IL2CPP \\
    Visual Studio Code & 1.143.1 & edytor kodu \\
    Blender & 5.1.2 & modele i eksport FBX \\
    Python & 3.12.10 & prototyp reguł i analiza wyników \\
    Git & 2.53.0 & kontrola wersji \\
    \bottomrule
  \end{tabular}
\end{table}

\CoTuMaByc{Zależności z plików \kol{*.asmdef}.}

\begin{table}
  \centering
  \footnotesize
  \caption{Biblioteki kodu generatora}
  \label{tab:biblioteki}
  \begin{tabularx}{\linewidth}{@{}llLl@{}}
    \toprule
    Biblioteka & Folder & Zależy od bibliotek projektu & Platforma \\
    \midrule
    \kol{WFC.Core} & \kol{Assets/Scripts/Core} & brak & wszystkie \\
    \kol{WFC.Runtime} & \kol{Assets/Scripts/Runtime} & \kol{WFC.Core} & wszystkie \\
    \kol{WFC.Editor} & \kol{Assets/Scripts/Editor} & \kol{WFC.Core}, \kol{WFC.Runtime} & edytor \\
    \kol{WFC.Tests} & \kol{Assets/Scripts/Tests} & \kol{WFC.Core}, \kol{WFC.Runtime}, \kol{WFC.Editor} & edytor \\
    \bottomrule
  \end{tabularx}
\end{table}

\CoTuMaByc{Liczby z sekcji kontrolnej \kol{Assets/Data/Rules/wfc_rules_zamek.json}.}

\begin{table}
  \centering
  \footnotesize
  \caption{Parametry zestawu kafelków}
  \label{tab:parametry-regul}
  \begin{tabular}{@{}lr@{}}
    \toprule
    Parametr & Wartość \\
    \midrule
    Moduł siatki & 3\,$\times$\,3\,$\times$\,3 m \\
    Modele & 43 \\
    Materiały & 12 \\
    Kafelki & 50 \\
    Stany (kafelek z obrotem) & 170 \\
    Obroty kafelka & 1 albo 4 \\
    Etykiety boczne & 36 \\
    Etykiety pionowe & 13 \\
    Pary zgodnych etykiet & 33 \\
    \bottomrule
  \end{tabular}
\end{table}

\begin{table}
  \centering
  \footnotesize
  \caption{Zestawy kafelków o rosnącej złożoności}
  \label{tab:zestawy}
  \begin{tabular}{@{}lrr@{}}
    \toprule
    Zestaw & Kafelki & Stany \\
    \midrule
    \kol{Prosty} & 8 & 26 \\
    \kol{Sredni} & 22 & 79 \\
    \kol{PelnyBezDekoracji} & 41 & 146 \\
    \kol{Pelny} & 50 & 170 \\
    \bottomrule
  \end{tabular}
\end{table}

\CoTuMaByc{Opisy z \kol{Assets/Scripts/Core/Rules/RuleVariant.cs}; numery zasad z arkusza reguł.}

\begin{table}
  \centering
  \footnotesize
  \caption{Warianty reguł powstałe z zestawu Pelny}
  \label{tab:warianty-regul}
  \begin{tabularx}{\linewidth}{@{}lL@{}}
    \toprule
    Wariant & Zmiana \\
    \midrule
    \kol{PelnyBezWolnegoPola} & bez zasad 12 i 23: boki murów oraz zewnętrzne ściany wież i donżonów mogą się stykać bezpośrednio \\
    \kol{PelnyBezDrzwiNadSchodami} & bez zasady 35: powietrze nad schodami może stać przed ślepą ścianą wieży albo donżonu i przed okapem dachu \\
    \kol{PelnyTrawa040} & wagi 5 kafelków trawy $\times$ 0,4 \\
    \kol{PelnyTrawa025} & wagi 5 kafelków trawy $\times$ 0,25 \\
    \bottomrule
  \end{tabularx}
\end{table}

\begin{table}
  \centering
  \footnotesize
  \caption{Rozwiązywacze}
  \label{tab:rozwiazywacze}
  \begin{tabular}{@{}lll@{}}
    \toprule
    Wartość w \kol{rozwiazywacz} & Przechowywanie możliwości & Cofanie i reset lokalny \\
    \midrule
    \kol{Naive} & listy & tak \\
    \kol{SupportCount} & liczniki wsparcia & nie \\
    \kol{Bitset} & maski bitowe & tak \\
    \kol{Burst} & maski bitowe, kompilator Burst & tak \\
    \bottomrule
  \end{tabular}
\end{table}

\begin{table}
  \centering
  \footnotesize
  \caption{Heurystyki wyboru komórki}
  \label{tab:heurystyki}
  \begin{tabular}{@{}ll@{}}
    \toprule
    Wartość w \kol{kolejnosc} & Kolejność \\
    \midrule
    \kol{LevelFirst} & piętrami, w piętrze najmniejsza entropia \\
    \kol{Scanline} & po kolei \\
    \kol{MinRemaining} & Minimum Remaining Values (MRV): najmniej możliwości \\
    \kol{Entropy} & najmniejsza entropia \\
    \kol{Random} & losowo \\
    \bottomrule
  \end{tabular}
\end{table}

\begin{table}
  \centering
  \footnotesize
  \caption{Reakcje na sprzeczność}
  \label{tab:reakcje}
  \begin{tabularx}{\linewidth}{@{}lL@{}}
    \toprule
    Wartość w \kol{sprzecznosc} & Działanie \\
    \midrule
    \kol{Restart} & start od nowa; limit prób na mapę 1000, w E9: 30 \\
    \kol{Backtrack} & cofanie decyzji, do 1000 cofnięć na próbę \\
    \kol{LocalReset} & reset sześcianu 5\,$\times$\,5\,$\times$\,5 komórek wokół sprzeczności, do 100 resetów na próbę \\
    \bottomrule
  \end{tabularx}
\end{table}

\begin{table}
  \centering
  \footnotesize
  \caption{Sposoby rysowania mapy}
  \label{tab:rysowanie}
  \begin{tabular}{@{}ll@{}}
    \toprule
    Wartość w \kol{rysowanie} & Sposób \\
    \midrule
    \kol{GameObjects} & obiekty w scenie \\
    \kol{Instanced} & instancje na karcie graficznej \\
    \kol{Combined} & scalone siatki \\
    \kol{CombinedAtlas} & scalone siatki z atlasem \\
    \bottomrule
  \end{tabular}
\end{table}

\CoTuMaByc{Ustawienia zasobu URP z \kol{Assets/Scripts/Runtime/Generation/RenderBatching.cs}.}

\begin{table}
  \centering
  \footnotesize
  \caption{Ustawienia łączenia rysowania w URP}
  \label{tab:laczenie}
  \begin{tabular}{@{}lcc@{}}
    \toprule
    Wartość w \kol{laczenie_rysowania} & SRP Batcher & GPU Resident Drawer \\
    \midrule
    \kol{SrpBatcherAndResidentDrawer} & tak & tak \\
    \kol{SrpBatcher} & tak & nie \\
    \kol{InstancingOnly} & nie & nie \\
    \bottomrule
  \end{tabular}
\end{table}

\begin{table}
  \centering
  \footnotesize
  \caption{Rozmiary siatki}
  \label{tab:rozmiary}
  \begin{tabular}{@{}lrl@{}}
    \toprule
    E\,$\times$\,N\,$\times$\,piętra & Komórki & Wymiar mapy [m] \\
    \midrule
    \siatka{10}{10}{8} & 800 & \siatka{30}{30}{24} \\
    \siatka{20}{20}{13} & 5200 & \siatka{60}{60}{39} \\
    \siatka{40}{40}{13} & 20\,800 & \siatka{120}{120}{39} \\
    \siatka{80}{80}{13} & 83\,200 & \siatka{240}{240}{39} \\
    \bottomrule
  \end{tabular}
\end{table}

\CoTuMaByc{Ustawienia z \kol{Assets/Scripts/Editor/ExperimentPresets.cs}. Bez innego zapisu: pierwsze ziarno 1, 1 powtórzenie, rozgrzewka 5 uruchomień na wariant, limit 1000 prób. Liczba ziaren w zasobach: E1--E6 i E8: 100, E7: 30, E9: 20. Argument \kol{-seeds} ją nadpisuje, więc wpisz tę, której naprawdę użyto.}

\begin{table}
  \centering
  \footnotesize
  \caption{Plan eksperymentów}
  \label{tab:eksperymenty}
  \begin{tabularx}{\linewidth}{@{}lLrrL@{}}
    \toprule
    Nr & Co się zmienia & Warianty & Ziarna & Ustawienia stałe \\
    \midrule
    E1 & rozwiązywacz (4) & 4 & 100 & \kol{Pelny}, \siatka{20}{20}{13}, \kol{LevelFirst}, \kol{Restart} \\
    E2 & heurystyka (5) & 5 & 100 & \kol{Pelny}, \siatka{20}{20}{13}, \kol{Bitset}, \kol{Restart} \\
    E3 & reakcja (3) $\times$ rozmiar (4) & 12 & 100 & \kol{Pelny}, \kol{Bitset}, \kol{LevelFirst}, limit 600 s na mapę \\
    E4 & rozwiązywacz (4) $\times$ rozmiar (4) & 16 & 100 & \kol{Pelny}, \kol{LevelFirst}, \kol{Restart}, limit 600 s na mapę \\
    E5 & zestaw (4) & 4 & 100 & \siatka{20}{20}{13}, \kol{Bitset}, \kol{LevelFirst}, \kol{Restart} \\
    E6 & \kol{Pelny} i warianty reguł (4) & 5 & 100 & \siatka{20}{20}{13}, \kol{Bitset}, \kol{LevelFirst}, \kol{Restart} \\
    E7 & rysowanie (4) $\times$ łączenie (3) $\times$ rozmiar (4) & 48 & 30 & \kol{Pelny}, \kol{Burst}, \kol{LevelFirst}, \kol{Backtrack}, rozgrzewka 3, 120 klatek na mapę, tylko w oknie \\
    E8 & 4 zestawy i 4 warianty reguł & 8 & 100 & \siatka{20}{20}{13}, \kol{Burst}, \kol{LevelFirst}, \kol{Restart} \\
    E9 & zaostrzenie reguł (9 poziomów) $\times$ 10 losowań & 81 & 20 & \kol{Pelny}, \siatka{20}{20}{13}, \kol{Burst}, \kol{LevelFirst}, \kol{Restart}, rozgrzewka 1, limit 30 prób \\
    \bottomrule
  \end{tabularx}
\end{table}

\CoTuMaByc{Część danych jest w każdym wierszu pliku wyników: \kol{komputer}, \kol{procesor}, \kol{ram_gb}, \kol{wersja_unity}, \kol{wersja_programu}, \kol{srodowisko}. Do pracy tylko wiersze z \kol{srodowisko} równym \kol{IL2CPP}. Karty graficznej plik wyników nie zapisuje: model i wersję sterownika trzeba zanotować z komputera pomiarowego. Rozdzielczość i rodzaj aplikacji wpisane według polecenia uruchomienia badań i ustawień \kol{AppBuilder}.}

\begin{table}
  \centering
  \footnotesize
  \caption{Środowisko pomiarowe}
  \label{tab:sprzet}
  \begin{tabular}{@{}ll@{}}
    \toprule
    Element & Wartość \\
    \midrule
    Procesor & AMD Ryzen 7 9800X3D, 8 rdzeni \\
    Pamięć operacyjna & 32 GB DDR5-6000, CL30 \\
    Karta graficzna & ASUS PRIME GeForce RTX 5080 OC Edition, 16 GB GDDR7 \\
    Sterownik karty graficznej & NVIDIA 610.88 \\
    System operacyjny & Windows 11 Pro \\
    Rozdzielczość okna & 1600\,$\times$\,900 pikseli, aplikacja w oknie \\
    Wersja Unity & 6000.6.2f1 \\
    Rodzaj aplikacji & IL2CPP, wersja Release \\
    \bottomrule
  \end{tabular}
\end{table}

\begin{table}
  \centering
  \footnotesize
  \caption{Główne metryki}
  \label{tab:metryki}
  \begin{tabular}{@{}lll@{}}
    \toprule
    Metryka & Kolumna pliku wyników & Jednostka \\
    \midrule
    Czas generacji & \kol{czas_ms} & ms \\
    Czas wyboru komórki & \kol{czas_wybor_ms} & ms \\
    Czas propagacji & \kol{czas_propagacja_ms} & ms \\
    Czas cofania & \kol{czas_cofanie_ms} & ms \\
    Próby i sprzeczności & \kol{proby}, \kol{sprzecznosci} & liczba \\
    Naruszenia zasad & \kol{naruszenia} & liczba \\
    Pamięć szacowana & \kol{pamiec_szczytowa_mb} & MB \\
    Pamięć rzeczywista & \kol{pamiec_rzeczywista_mb} & MB \\
    Wywołania rysowania & \kol{wywolania_rysowania} & liczba \\
    Narysowane obiekty & \kol{narysowane_obiekty} & liczba \\
    Zmiany stanu SetPass & \kol{setpass} & liczba \\
    Czas klatki & \kol{czas_klatki_ms} & ms \\
    Najdłuższa klatka & \kol{najdluzsza_klatka_ms} & ms \\
    Czas stawiania mapy & \kol{czas_stawiania_ms} & ms \\
    Pamięć silnika & \kol{pamiec_unity_mb} & MB \\
    \bottomrule
  \end{tabular}
\end{table}

\CoTuMaByc{Znaczniki z \kol{Assets/Scripts/Core/Generation/WfcMarkers.cs}. Czasy w plikach wyników liczy osobno \kol{Stopwatch}.}

\begin{table}
  \centering
  \footnotesize
  \caption{Znaczniki w oknie Unity Profiler}
  \label{tab:znaczniki}
  \begin{tabular}{@{}ll@{}}
    \toprule
    Znacznik & Etap \\
    \midrule
    \kol{WFC.Generacja} & cała generacja mapy \\
    \kol{WFC.WyborKomorki} & wybór komórki \\
    \kol{WFC.Propagacja} & propagacja ograniczeń \\
    \kol{WFC.Cofanie} & cofanie decyzji \\
    \kol{WFC.SprawdzenieMapy} & sprawdzenie gotowej mapy \\
    \kol{WFC.StawianieMapy} & stawianie mapy \\
    \kol{WFC.Scalanie} & scalanie siatek \\
    \bottomrule
  \end{tabular}
\end{table}

\CoTuMaByc{Testy z \kol{Assets/Scripts/Tests}. Liczby jak w oknie Test Runner: test z \kol{TestCase} albo \kol{ValueSource} liczy się raz na każdy zestaw argumentów. Ostatni przebieg: 4.11.2026, 280 testów, wszystkie zaliczone.}

\begin{table}
  \centering
  \footnotesize
  \caption{Testy automatyczne}
  \label{tab:testy}
  \begin{tabularx}{\linewidth}{@{}lLr@{}}
    \toprule
    Obszar & Pliki testów & Testy \\
    \midrule
    Reguły i kafelki & \kol{RuleCompilerTests}, \kol{RulesControlTests}, \kol{RuleVariantTests}, \kol{TilesetValidatorTests}, \kol{OrientationTests}, \kol{DirectionTests}, \kol{SpreadsheetTests}, \kol{ImportedAssetsTests} & 115 \\
    Generacja & \kol{GenerationTests}, \kol{SolverTests}, \kol{SelectorTests}, \kol{ContradictionTests}, \kol{StepPreviewTests} & 80 \\
    Analiza map & \kol{RuleCheckerTests}, \kol{StructureAnalyzerTests}, \kol{DiversityTests} & 18 \\
    Rysowanie i narzędzie & \kol{MapBuilderTests}, \kol{CombinedMapTests}, \kol{GeneratorWindowTests}, \kol{FlyCameraTests} & 29 \\
    Pomiary & \kol{BenchmarkTests}, \kol{MeasurementTests} & 38 \\
    \midrule
    Razem & 22 pliki & 280 \\
    \bottomrule
  \end{tabularx}
\end{table}

\CoTuMaByc{Pusty układ do analizy wymagań; treść i liczbę wierszy ustalasz we własnym zakresie.}

\begin{table}
  \centering
  \footnotesize
  \caption{Wymagania}
  \label{tab:wymagania}
  \begin{tabularx}{\linewidth}{@{}lLL@{}}
    \toprule
    Oznaczenie & Wymaganie & Sposób sprawdzenia \\
    \midrule
    \wyn & \wyn & \wyn \\
    \wyn & \wyn & \wyn \\
    \wyn & \wyn & \wyn \\
    \bottomrule
  \end{tabularx}
\end{table}

\CoTuMaByc{Wymaganie Wydziału: umiejscowienie zagadnienia na tle podobnych rozwiązań. Wiersz „Ta praca” ma dane projektu; pozostałe wiersze to propozycje do sprawdzenia, a kolumny i wnioski ustalasz we własnym zakresie.}

\begin{table}
  \centering
  \footnotesize
  \caption{Podobne rozwiązania}
  \label{tab:podobne}
  \begin{tabular}{@{}lccc@{}}
    \toprule
    Rozwiązanie & Język lub silnik & Siatka 3D & Rysowanie \\
    \midrule
    Gumin \cite{Gumin2016} & \wyn & \wyn & \wyn \\
    DeBroglie \cite{DeBroglie} & \wyn & \wyn & \wyn \\
    marian42 \cite{Marian42WFC} & \wyn & \wyn & \wyn \\
    Ta praca & C\#, Unity 6 & tak & 4 sposoby \\
    \bottomrule
  \end{tabular}
\end{table}


\section{Wzory}
%%% KLOCKI: WZORY
%%% Wzory zapisane według tego, co liczy kod. Oznaczenia i opis dopisujesz pod wzorami.
%%% Numeruj tylko te równania, do których się odwołujesz.

\CoTuMaByc{Entropia komórki, kod: \kol{Assets/Scripts/Core/Generation/WeightedChooser.cs}. Źródło pojęcia: \cite{Shannon1948}. $D(c)$: możliwe stany komórki $c$, $w_s$: waga stanu $s$.}

\begin{equation}
  \begin{gathered}
    H(c) = -\sum_{s \in D(c)} p_s \ln p_s
         = \ln W_c - \frac{1}{W_c} \sum_{s \in D(c)} w_s \ln w_s , \\
    p_s = \frac{w_s}{W_c} , \qquad W_c = \sum_{s \in D(c)} w_s
  \end{gathered}
  \label{eq:entropia}
\end{equation}

\CoTuMaByc{Trzy miary restrykcyjności, kod: \kol{Assets/Scripts/Core/Analysis/RuleRestrictiveness.cs}. $m$: liczba stanów, $P$: liczba dozwolonych trójek (stan, kierunek, sąsiad), $N_k(s)$: dozwoleni sąsiedzi stanu $s$ w kierunku $k$, $W$: suma wag. Kolumny: \kol{gestosc_regul_proc}, \kol{gestosc_wazona_proc}, \kol{kappa}. Współczynnik $\kappa$: \cite{Gent1996}; w siatce 3D przyjęto $c/n = 3$.}

\begin{equation}
  d = \frac{P}{6 m^2}
  \label{eq:gestosc}
\end{equation}

\begin{equation}
  d_w = \frac{1}{6} \sum_{k=1}^{6} \sum_{s=1}^{m} \frac{w_s}{W} \sum_{t \in N_k(s)} \frac{w_t}{W}
  \label{eq:gestosc-wazona}
\end{equation}

\begin{equation}
  \kappa = -\frac{c}{n} \log_m d
  \label{eq:kappa}
\end{equation}

\CoTuMaByc{Odsetek udanych prób, kolumna \kol{udane_proby_proc} w pliku \kol{*_warianty.csv}. $M$: udane mapy, $A$: wszystkie próby.}

\begin{equation}
  u = 100 \cdot \frac{M}{A}
  \label{eq:udane-proby}
\end{equation}

\CoTuMaByc{Udział propagacji w usunięciach i usunięcia na decyzję. Kolumny: \kol{usuniete_propagacja} ($r_p$), \kol{usuniete_decyzje} ($r_d$), \kol{decyzje} ($D$).}

\begin{equation}
  q = \frac{r_p}{r_p + r_d} , \qquad e = \frac{r_p}{D}
  \label{eq:propagacja}
\end{equation}

\CoTuMaByc{Liczba obiektów na wywołanie rysowania. Kolumny: \kol{narysowane_obiekty} ($O$), \kol{wywolania_rysowania} ($V$).}

\begin{equation}
  b = \frac{O}{V}
  \label{eq:obiekty-na-wywolanie}
\end{equation}

\CoTuMaByc{Rozstęp międzykwartylowy oraz korelacja rang: \cite{Spearman1904}. Kod: \kol{Assets/Scripts/Core/Analysis/Statistics.cs}.}

\begin{equation}
  \mathrm{IQR} = Q_3 - Q_1
  \label{eq:iqr}
\end{equation}

\begin{equation}
  \rho = \mathrm{corr}\bigl(\mathrm{rg}(X), \mathrm{rg}(Y)\bigr)
  \label{eq:spearman}
\end{equation}

\CoTuMaByc{Nachylenie prostej w skali logarytmicznej dla czasu lub pamięci $t$ wobec liczby komórek $n$ (\cref{tab:e4-czas,tab:e4-pamiec}).}

\begin{equation}
  \log t = a \log n + b
  \label{eq:nachylenie}
\end{equation}


\section{Pseudokody}
%%% KLOCKI: MIEJSCA NA PSEUDOKOD
%%% Kroki wpisujesz w miejsce pola w nawiasie. Składnia: \State, \For{...} \EndFor, \If{...} \EndIf, \While{...} \EndWhile.
%%% Podpisy są robocze, a pole mówi tylko, gdzie w kodzie jest dana część.

\begin{algorithm}
  \caption{Główna pętla generacji}
  \label{alg:petla}
  \begin{algorithmic}[1]
    \State \Wpisz{kroki; kod: \kol{SolverBase.cs}}
  \end{algorithmic}
\end{algorithm}

\begin{algorithm}
  \caption{Propagacja ograniczeń na listach}
  \label{alg:propagacja-listy}
  \begin{algorithmic}[1]
    \State \Wpisz{kroki; kod: \kol{NaiveSolver.cs}, \kol{CellQueue.cs}}
  \end{algorithmic}
\end{algorithm}

\begin{algorithm}
  \caption{Propagacja z licznikami wsparcia}
  \label{alg:propagacja-liczniki}
  \begin{algorithmic}[1]
    \State \Wpisz{kroki; kod: \kol{SupportCountSolver.cs}}
  \end{algorithmic}
\end{algorithm}

\begin{algorithm}
  \caption{Propagacja na maskach bitowych}
  \label{alg:propagacja-maski}
  \begin{algorithmic}[1]
    \State \Wpisz{kroki; kod: \kol{BitsetPropagation.cs}}
  \end{algorithmic}
\end{algorithm}

\begin{algorithm}
  \caption{Wybór komórki}
  \label{alg:wybor-komorki}
  \begin{algorithmic}[1]
    \State \Wpisz{kroki; kod: \kol{CellSelectors.cs}, \kol{CellPriorityQueue.cs}}
  \end{algorithmic}
\end{algorithm}

\begin{algorithm}
  \caption{Cofanie decyzji}
  \label{alg:cofanie}
  \begin{algorithmic}[1]
    \State \Wpisz{kroki; kod: \kol{IContradictionHandler.cs}}
  \end{algorithmic}
\end{algorithm}

\begin{algorithm}
  \caption{Reset lokalny}
  \label{alg:reset-lokalny}
  \begin{algorithmic}[1]
    \State \Wpisz{kroki; kod: \kol{IContradictionHandler.cs}}
  \end{algorithmic}
\end{algorithm}

\begin{algorithm}
  \caption{Scalanie siatek}
  \label{alg:scalanie}
  \begin{algorithmic}[1]
    \State \Wpisz{kroki; kod: \kol{CombinedMap.cs}}
  \end{algorithmic}
\end{algorithm}


\section{Listingi}
%%% KLOCKI: MIEJSCA NA LISTINGI
%%% Wybrany fragment kodu zapisz w folderze listingi pod podaną nazwą, a listing pojawi się sam.
%%% Pliki w UTF-8 bez znacznika BOM. Fragmenty powstałe z pomocą AI oznacz w komentarzu kodu.
%%% Stan 12.11.2026: wszystkie 10 plików wycięte z kodu (commit 4e10be9) bez komentarzy ///, z wcięciem po 2 spacje; opuszczenia oznacza // ...;
%%% pierwsza linia każdego pliku to oznaczenie fragmentu z AI.

\MiejsceNaListing{kompilacja-regul}{Kompilacja reguł}{lst:kompilacja-regul}{Kod: \kol{Assets/Scripts/Core/Rules/RuleCompiler.cs}.}

\MiejsceNaListing{maska-propagacja}{Propagacja na maskach bitowych}{lst:maska-propagacja}{Kod: \kol{Assets/Scripts/Core/Generation/BitsetPropagation.cs}.}

\MiejsceNaListing{zadanie-burst}{Rozwiązywacz w kompilatorze Burst}{lst:zadanie-burst}{Kod: \kol{Assets/Scripts/Core/Generation/BurstSolver.cs}.}

\MiejsceNaListing{ziarno}{Generator liczb losowych}{lst:ziarno}{Kod: \kol{Assets/Scripts/Core/Generation/DeterministicRandom.cs}.}

\MiejsceNaListing{uklad-wierzcholka}{Układ wierzchołka scalonej siatki}{lst:uklad-wierzcholka}{Kod: pole \kol{VertexLayout} w \kol{Assets/Scripts/Runtime/Generation/CombinedMap.cs}.}

\MiejsceNaListing{zadanie-fragmentu}{Zadanie zapisujące fragment siatki}{lst:zadanie-fragmentu}{Kod: struktura \kol{WriteChunkJob} w \kol{CombinedMap.cs}.}

\MiejsceNaListing{zapis-siatek}{Zapis siatek przez Mesh API}{lst:zapis-siatek}{Kod: \kol{Mesh.AllocateWritableMeshData}, \kol{SetVertexBufferParams} i \kol{Mesh.ApplyAndDisposeWritableMeshData} w \kol{CombinedMap.cs}. Dokumentacja: \cite{UnityAllocateWritableMeshData,UnitySetVertexBufferParams}.}

\MiejsceNaListing{zwalnianie-pamieci}{Zwalnianie pamięci mapy}{lst:zwalnianie-pamieci}{Kod: metoda \kol{Dispose} w \kol{CombinedMap.cs}.}

\MiejsceNaListing{atlas}{Budowa atlasu materiałów}{lst:atlas}{Kod: \kol{Assets/Scripts/Runtime/Generation/MaterialAtlas.cs}.}

\MiejsceNaListing{znaczniki}{Znaczniki dla okna Profiler}{lst:znaczniki}{Kod: \kol{Assets/Scripts/Core/Generation/WfcMarkers.cs}. Dokumentacja: \cite{UnityProfilerMarker}.}


\section{Rysunki}
%%% KLOCKI: MIEJSCA NA RYSUNKI
%%% Plik zapisz w folderze figures pod podaną nazwą (pdf, png albo jpg), a rysunek pojawi się sam.
%%% Podpisy są robocze. Wykresy tworzy skrypt Analiza/wykresy.py od razu w folderze figures (python Analiza/wykresy.py);
%%% szerokość pliku jest równa szerokości w nawiasie kwadratowym, więc nie zmieniaj jej bez ponownego uruchomienia skryptu.
%%% Zrzuty ekranu przygotowujesz we własnym zakresie. Opis pokazuje tylko, z czego powstaje rysunek.

\MiejsceNaRysunek{architektura}{Zależności bibliotek kodu}{fig:architektura}{Schemat 4 bibliotek, \cref{tab:biblioteki}.}

\MiejsceNaRysunek{modele}{Modele zestawu kafelków}{fig:modele}{Podgląd 43 modeli; scena z \kol{TilesetPreviewSceneBuilder}.}

\MiejsceNaRysunek{orientacja}{Etykiety ścian i obroty kafelka}{fig:orientacja}{Przykład: \kol{Mur_Dol} w 4 obrotach.}

\begin{figure}
  \centering
  \begin{subfigure}[t]{0.32\linewidth}
    \centering
    \MiejsceNaObraz[\linewidth]{krok-1}{Zrzut z podglądu krok po kroku}
    \caption{Etap pierwszy}
  \end{subfigure}\hfill
  \begin{subfigure}[t]{0.32\linewidth}
    \centering
    \MiejsceNaObraz[\linewidth]{krok-2}{Zrzut z podglądu krok po kroku}
    \caption{Etap drugi}
  \end{subfigure}\hfill
  \begin{subfigure}[t]{0.32\linewidth}
    \centering
    \MiejsceNaObraz[\linewidth]{krok-3}{Zrzut z podglądu krok po kroku}
    \caption{Etap trzeci}
  \end{subfigure}
  \caption{Generacja mapy krok po kroku}
  \label{fig:krok-po-kroku}
\end{figure}

\MiejsceNaRysunek{okno-generatora}{Okno Generator zamków w edytorze}{fig:okno-generatora}{Zrzut okna z \kol{GeneratorWindow}.}

\MiejsceNaRysunek{prefab}{Mapa zapisana jako prefab}{fig:prefab}{Zrzut prefabu z \kol{MapBaker}.}

\MiejsceNaRysunek{aplikacja}{Aplikacja podczas badania}{fig:aplikacja}{Zrzut okna z panelem badań i zamkiem.}

\MiejsceNaRysunek[0.5\linewidth]{atlas}{Atlas materiałów}{fig:atlas}{Tekstura z \kol{MaterialAtlas}.}

\MiejsceNaRysunek{profiler}{Znaczniki w oknie Unity Profiler}{fig:profiler}{Zrzut ze znacznikami \kol{WFC.*}, \cref{tab:znaczniki}.}

\MiejsceNaRysunek{e1-czas}{E1: czas generacji według rozwiązywacza}{fig:e1-czas}{Wykres pudełkowy, kolumna \kol{czas_ms}: pudełko od Q1 do Q3, wąsy do 1,5 rozstępu międzykwartylowego, kropki to pojedyncze mapy.}

\MiejsceNaRysunek[\linewidth]{e2-czas-udane}{E2: czas generacji i odsetek udanych prób według heurystyki}{fig:e2-czas-udane}{Panel a: wykres pudełkowy kolumny \kol{czas_ms}; panel b: \kol{udane_proby_proc} liczone jak w \kol{*_warianty.csv}.}

\begin{figure}
  \centering
  \begin{subfigure}[t]{0.32\linewidth}
    \centering
    \MiejsceNaObraz[\linewidth]{e2-levelfirst}{Zrzut zamku}
    \caption{LevelFirst}
  \end{subfigure}\hfill
  \begin{subfigure}[t]{0.32\linewidth}
    \centering
    \MiejsceNaObraz[\linewidth]{e2-scanline}{Zrzut zamku}
    \caption{Scanline}
  \end{subfigure}\hfill
  \begin{subfigure}[t]{0.32\linewidth}
    \centering
    \MiejsceNaObraz[\linewidth]{e2-minremaining}{Zrzut zamku}
    \caption{MinRemaining}
  \end{subfigure}

  \medskip
  \begin{subfigure}[t]{0.32\linewidth}
    \centering
    \MiejsceNaObraz[\linewidth]{e2-entropy}{Zrzut zamku}
    \caption{Entropy}
  \end{subfigure}\hspace{0.02\linewidth}
  \begin{subfigure}[t]{0.32\linewidth}
    \centering
    \MiejsceNaObraz[\linewidth]{e2-random}{Zrzut zamku}
    \caption{Random}
  \end{subfigure}
  \caption{E2: zamki z tego samego ziarna dla pięciu heurystyk}
  \label{fig:e2-zamki}
\end{figure}

\MiejsceNaRysunek{e3-czas-rozmiar}{E3: czas generacji wobec rozmiaru siatki}{fig:e3-czas-rozmiar}{Mediana i rozstęp międzykwartylowy kolumny \kol{czas_ms} według \kol{sprzecznosc}; obie osie w skali logarytmicznej. Linia limitu 600 s pojawia się, gdy wyniki do niego dochodzą.}

\MiejsceNaRysunek{e4-czas-log}{E4: czas generacji wobec liczby komórek}{fig:e4-czas-log}{Mediana i rozstęp międzykwartylowy kolumny \kol{czas_ms} według \kol{rozwiazywacz}; obie osie w skali logarytmicznej. Nachylenie $a$ (\cref{eq:nachylenie}) skrypt wypisuje w konsoli.}

\MiejsceNaRysunek{e4-pamiec-log}{E4: pamięć wobec liczby komórek}{fig:e4-pamiec-log}{Jak wyżej dla \kol{pamiec_rzeczywista_mb}. Mediana równa 0 nie mieści się na skali logarytmicznej; skrypt wypisuje wtedy, czego brakuje.}

\MiejsceNaRysunek{e5-czas-stany}{E5: czas generacji wobec liczby stanów}{fig:e5-czas-stany}{Mediana i rozstęp międzykwartylowy kolumny \kol{czas_ms} według \kol{zestaw}; stany z \cref{tab:zestawy}.}

\MiejsceNaRysunek[\linewidth]{e6-udane-zabudowa}{E6: odsetek udanych prób i zabudowa według wariantu reguł}{fig:e6-udane-zabudowa}{Kolumny \kol{udane_proby_proc} i \kol{zabudowa_proc} liczone jak w \kol{*_warianty.csv}.}

\MiejsceNaRysunek[\linewidth]{e7-wywolania}{E7: wywołania rysowania wobec rozmiaru siatki}{fig:e7-wywolania}{Mediana i rozstęp międzykwartylowy kolumny \kol{wywolania_rysowania}; panele według \kol{laczenie_rysowania}, linie według \kol{rysowanie}.}

\MiejsceNaRysunek[\linewidth]{e7-klatka}{E7: czas klatki wobec rozmiaru siatki}{fig:e7-klatka}{Jak wyżej dla \kol{czas_klatki_ms}.}

\MiejsceNaRysunek[\linewidth]{e8-restrykcyjnosc}{E8: restrykcyjność reguł a odsetek udanych prób}{fig:e8-restrykcyjnosc}{Wykres punktowy: \kol{gestosc_regul_proc}, \kol{gestosc_wazona_proc} i \kol{kappa} wobec \kol{udane_proby_proc}; pusty znacznik to mapy zdegenerowane.}

\MiejsceNaRysunek{e9-zaostrzanie}{E9: odsetek udanych prób wobec zaostrzenia reguł}{fig:e9-zaostrzanie}{Kolumny \kol{udane_proby_proc} i \kol{zabudowa_proc} według \kol{zaostrzenie_proc}; kropka to jedno losowanie reguł, linia to mediana losowań.}

\MiejsceNaRysunek{propagacja-udzial}{Udział propagacji w usuniętych możliwościach}{fig:propagacja-udzial}{Mapy rozwiązywacza \kol{Bitset} z E4; kolumny \kol{usuniete_propagacja}, \kol{usuniete_decyzje}, \cref{eq:propagacja}.}


\section{Wyniki pomiarów}
%%% KLOCKI: TABELE WYNIKÓW
%%% Wiersze i kolumny wynikają z ustawień eksperymentów. Wyniki wpisujesz w miejsce \wyn.
%%% Do pracy tylko wyniki z aplikacji IL2CPP w wersji Release (kolumna srodowisko równa IL2CPP).
%%% Pliki z folderu Wyniki: główny E1_*.csv oraz podsumowania E1_*_warianty.csv i E1_*_korelacje.csv.

\CoTuMaByc{Z \kol{E1_*_warianty.csv}: \kol{czas_ms_mediana}, \kol{czas_ms_q1}, \kol{czas_ms_q3}, \kol{pamiec_rzeczywista_mb_mediana}. Z głównego pliku: \kol{czas_propagacja_ms}.}

\begin{table}
  \centering
  \footnotesize
  \caption{E1: czas generacji i pamięć według rozwiązywacza}
  \label{tab:e1}
  \begin{tabularx}{\linewidth}{@{}lCCCCC@{}}
    \toprule
    Rozwiązywacz & Czas: mediana [ms] & Czas: Q1 [ms] & Czas: Q3 [ms] & Propagacja: mediana [ms] & Pamięć rzeczywista [MB] \\
    \midrule
    \kol{Naive} & 837,8 & 803,4 & 1449,1 & 825,2 & 6,12 \\
    \kol{SupportCount} & 81,3 & 79,9 & 132,0 & 67,3 & 13,43 \\
    \kol{Bitset} & 26,3 & 24,7 & 36,4 & 12,1 & 1,92 \\
    \kol{Burst} & 20,1 & 18,4 & 25,3 & 5,8 & 1,92 \\
    \bottomrule
  \end{tabularx}
\end{table}

\CoTuMaByc{Z \kol{*_mann_whitney.csv} ze skryptu \kol{Analiza/analiza_wynikow.py}: \kol{U}, \kol{p}. Test: \cite{MannWhitney1947}.}

\begin{table}
  \centering
  \footnotesize
  \caption{E1: test Manna-Whitneya dla par rozwiązywaczy}
  \label{tab:e1-mann-whitney}
  \begin{tabular}{@{}llrr@{}}
    \toprule
    Wariant A & Wariant B & $U$ & $p$ \\
    \midrule
    \kol{Naive} & \kol{SupportCount} & 10\,000 & $2{,}56 \cdot 10^{-34}$ \\
    \kol{Naive} & \kol{Bitset} & 10\,000 & $2{,}56 \cdot 10^{-34}$ \\
    \kol{Naive} & \kol{Burst} & 10\,000 & $2{,}56 \cdot 10^{-34}$ \\
    \kol{SupportCount} & \kol{Bitset} & 9928 & $2{,}19 \cdot 10^{-33}$ \\
    \kol{SupportCount} & \kol{Burst} & 10\,000 & $2{,}56 \cdot 10^{-34}$ \\
    \kol{Bitset} & \kol{Burst} & 8107 & $3{,}19 \cdot 10^{-14}$ \\
    \bottomrule
  \end{tabular}
\end{table}

\CoTuMaByc{Z \kol{E2_*_warianty.csv}: \kol{czas_ms_mediana}, \kol{proby}, \kol{sprzecznosci}, \kol{udane_proby_proc}.}

\begin{table}
  \centering
  \footnotesize
  \caption[E2: czas i sprzeczności według heurystyki]{E2: czas i sprzeczności według heurystyki. Próby i sprzeczności: suma ze 100 map}
  \label{tab:e2}
  \begin{tabularx}{\linewidth}{@{}lCCCC@{}}
    \toprule
    Heurystyka & Czas: mediana [ms] & Próby & Sprzeczności & Udane próby [\%] \\
    \midrule
    \kol{LevelFirst} & 26,4 & 139 & 39 & 71,9 \\
    \kol{Scanline} & 16,3 & 159 & 59 & 62,9 \\
    \kol{MinRemaining} & 226,7 & 2649 & 2549 & 3,8 \\
    \kol{Entropy} & 330,7 & 4086 & 3986 & 2,4 \\
    \kol{Random} & 381,7 & 5765 & 5665 & 1,7 \\
    \bottomrule
  \end{tabularx}
\end{table}

\CoTuMaByc{Z głównego pliku E2: \kol{wieze}, \kol{donzony}, \kol{mosty}, \kol{schody}, \kol{wys_wiez_sr}, \kol{wys_murow_sr}, \kol{trawa_dostepna_proc}.}

\begin{table}
  \centering
  \footnotesize
  \caption[E2: budowa map według heurystyki]{E2: budowa map według heurystyki, średnie na mapę}
  \label{tab:e2-budowa}
  \begin{tabularx}{\linewidth}{@{}lCCCCCCC@{}}
    \toprule
    Heurystyka & Wieże & Donżony & Mosty & Schody & Wysokość wież & Wysokość murów & Teren osiągalny [\%] \\
    \midrule
    \kol{LevelFirst} & 11,9 & 6,5 & 1,2 & 4,4 & 5,7 & 2,9 & 97,1 \\
    \kol{Scanline} & 16,6 & 4,5 & 1,5 & 3,7 & 5,9 & 2,9 & 94,1 \\
    \kol{MinRemaining} & 26,8 & 2,5 & 35,2 & 15,1 & 11,7 & 8,4 & 42,8 \\
    \kol{Entropy} & 29,2 & 1,7 & 26,6 & 10,9 & 11,6 & 9,0 & 41,6 \\
    \kol{Random} & 22,4 & 7,4 & 12,8 & 12,2 & 10,5 & 7,0 & 72,9 \\
    \bottomrule
  \end{tabularx}
\end{table}

\CoTuMaByc{Odsetek komórek innych niż w mapie poprzedniego ziarna. Z głównego pliku E2: \kol{roznica_mapy_proc}, \kol{roznica_parter_proc}, \kol{roznica_zabudowa_proc}.}

\begin{table}
  \centering
  \footnotesize
  \caption[E2: różnorodność map według heurystyki]{E2: różnorodność map według heurystyki, średni odsetek komórek innych niż w mapie poprzedniego ziarna}
  \label{tab:roznorodnosc}
  \begin{tabularx}{\linewidth}{@{}lCCC@{}}
    \toprule
    Heurystyka & Cała mapa [\%] & Parter [\%] & Zabudowa [\%] \\
    \midrule
    \kol{LevelFirst} & 23,9 & 79,5 & 98,2 \\
    \kol{Scanline} & 25,4 & 80,5 & 97,9 \\
    \kol{MinRemaining} & 67,8 & 82,7 & 94,6 \\
    \kol{Entropy} & 69,2 & 81,4 & 93,7 \\
    \kol{Random} & 58,5 & 83,2 & 96,2 \\
    \bottomrule
  \end{tabularx}
\end{table}

\CoTuMaByc{Z \kol{E3_*_warianty.csv}: \kol{czas_ms_mediana}. Limit 600 s na mapę; przekroczenie zapisuje się jako wynik.}

\begin{table}
  \centering
  \footnotesize
  \caption{E3: mediana czasu generacji [ms]}
  \label{tab:e3-czas}
  \begin{tabular}{@{}lrrrr@{}}
    \toprule
    Reakcja & 10$\times$10$\times$8 & 20$\times$20$\times$13 & 40$\times$40$\times$13 & 80$\times$80$\times$13 \\
    \midrule
    \kol{Restart} & 3,0 & 26,3 & 300,1 & 114\,750 \\
    \kol{Backtrack} & 3,0 & 25,7 & 117,1 & 496,6 \\
    \kol{LocalReset} & 3,0 & 25,4 & 116,7 & 497,2 \\
    \bottomrule
  \end{tabular}
\end{table}

\CoTuMaByc{Z \kol{E3_*_warianty.csv}: \kol{udane_mapy}, \kol{mapy}.}

\begin{table}
  \centering
  \footnotesize
  \caption{E3: udane mapy na liczbę map}
  \label{tab:e3-powodzenie}
  \begin{tabular}{@{}lrrrr@{}}
    \toprule
    Reakcja & 10$\times$10$\times$8 & 20$\times$20$\times$13 & 40$\times$40$\times$13 & 80$\times$80$\times$13 \\
    \midrule
    \kol{Restart} & 100/100 & 100/100 & 100/100 & 50/100 \\
    \kol{Backtrack} & 100/100 & 100/100 & 100/100 & 100/100 \\
    \kol{LocalReset} & 100/100 & 100/100 & 100/100 & 100/100 \\
    \bottomrule
  \end{tabular}
\end{table}

\CoTuMaByc{Z \kol{E4_*_warianty.csv}: \kol{czas_ms_mediana}. Nachylenie: \cref{eq:nachylenie}. Limit 600 s na mapę.}

\begin{table}
  \centering
  \footnotesize
  \caption[E4: mediana czasu generacji]{E4: mediana czasu generacji [ms]. Nachylenie $a$ dopasowano do trzech mniejszych siatek, bo przy 80$\times$80$\times$13 część map nie powstaje}
  \label{tab:e4-czas}
  \begin{tabular}{@{}lrrrrr@{}}
    \toprule
    Rozwiązywacz & 10$\times$10$\times$8 & 20$\times$20$\times$13 & 40$\times$40$\times$13 & 80$\times$80$\times$13 & Nachylenie $a$ \\
    \midrule
    Komórki & 800 & 5200 & 20\,800 & 83\,200 & -- \\
    \midrule
    \kol{Naive} & 90,0 & 835,5 & 12\,666 & 601\,188 & 1,50 \\
    \kol{SupportCount} & 10,3 & 82,9 & 1229,5 & 600\,000 & 1,45 \\
    \kol{Bitset} & 3,1 & 26,0 & 297,2 & 115\,059 & 1,39 \\
    \kol{Burst} & 2,3 & 19,8 & 207,3 & 73\,773 & 1,36 \\
    \bottomrule
  \end{tabular}
\end{table}

\CoTuMaByc{Z \kol{E4_*_warianty.csv}: \kol{pamiec_rzeczywista_mb_mediana}. Pojedyncze wartości pamięci szumią, więc porównuj mediany.}

\begin{table}
  \centering
  \footnotesize
  \caption{E4: mediana pamięci rzeczywistej [MB]}
  \label{tab:e4-pamiec}
  \begin{tabular}{@{}lrrrrr@{}}
    \toprule
    Rozwiązywacz & 10$\times$10$\times$8 & 20$\times$20$\times$13 & 40$\times$40$\times$13 & 80$\times$80$\times$13 & Nachylenie $a$ \\
    \midrule
    Komórki & 800 & 5200 & 20\,800 & 83\,200 & -- \\
    \midrule
    \kol{Naive} & 0,88 & 6,10 & 24,41 & 85,72 & 0,99 \\
    \kol{SupportCount} & 2,15 & 13,43 & 53,23 & 212,45 & 0,99 \\
    \kol{Bitset} & 0,26 & 1,92 & 7,57 & 30,22 & 1,02 \\
    \kol{Burst} & 0,26 & 1,92 & 7,57 & 30,22 & 1,02 \\
    \bottomrule
  \end{tabular}
\end{table}

\CoTuMaByc{Z \kol{E5_*_warianty.csv}: \kol{czas_ms_mediana}, \kol{czas_ms_q1}, \kol{czas_ms_q3}, \kol{sprzecznosci}.}

\begin{table}
  \centering
  \footnotesize
  \caption[E5: czas generacji według zestawu]{E5: czas generacji według zestawu. Sprzeczności: suma ze 100 map}
  \label{tab:e5}
  \setlength{\tabcolsep}{3pt}
  \begin{tabularx}{\linewidth}{@{}lCCCCCC@{}}
    \toprule
    Zestaw & Kafelki & Stany & Czas: mediana [ms] & Czas: Q1 [ms] & Czas: Q3 [ms] & Sprzeczności \\
    \midrule
    \kol{Prosty} & 8 & 26 & 5,8 & 4,4 & 8,2 & 117 \\
    \kol{Sredni} & 22 & 79 & 11,7 & 11,4 & 12,6 & 26 \\
    \kol{PelnyBezDekoracji} & 41 & 146 & 25,9 & 24,7 & 36,6 & 60 \\
    \kol{Pelny} & 50 & 170 & 26,0 & 24,7 & 36,4 & 39 \\
    \bottomrule
  \end{tabularx}
\end{table}

\CoTuMaByc{Z \kol{E6_*_warianty.csv}: \kol{sprzecznosci}, \kol{udane_proby_proc}, \kol{zabudowa_proc}. Z głównego pliku: \kol{wieze}.}

\begin{table}
  \centering
  \footnotesize
  \caption[E6: wyniki dla zestawu Pelny i wariantów reguł]{E6: wyniki dla zestawu Pelny i wariantów reguł. Sprzeczności: suma ze 100 map, wieże: średnia na mapę}
  \label{tab:e6}
  \begin{tabularx}{\linewidth}{@{}lCCCC@{}}
    \toprule
    Wariant & Sprzeczności & Udane próby [\%] & Zabudowa [\%] & Wieże \\
    \midrule
    \kol{Pelny} & 39 & 71,9 & 39,4 & 11,9 \\
    \kol{PelnyBezWolnegoPola} & 958 & 9,5 & 56,6 & 24,2 \\
    \kol{PelnyBezDrzwiNadSchodami} & 13 & 88,5 & 40,4 & 12,6 \\
    \kol{PelnyTrawa040} & 113 & 46,9 & 50,6 & 15,9 \\
    \kol{PelnyTrawa025} & 210 & 32,3 & 55,1 & 16,4 \\
    \bottomrule
  \end{tabularx}
\end{table}

\CoTuMaByc{Z \kol{E7_*_warianty.csv} albo mediana z głównego pliku: \kol{wywolania_rysowania_mediana}.}

\begin{table}
  \centering
  \footnotesize
  \caption{E7: mediana wywołań rysowania na klatkę}
  \label{tab:e7-wywolania}
  \begin{tabular}{@{}llrrrr@{}}
    \toprule
    Rysowanie & Łączenie rysowania & 10$\times$10$\times$8 & 20$\times$20$\times$13 & 40$\times$40$\times$13 & 80$\times$80$\times$13 \\
    \midrule
    \multirow{3}{*}{\kol{GameObjects}} & SRP Batcher + Resident Drawer & 115,5 & 183 & 211 & 211 \\
     & Sam SRP Batcher & 692 & 2916 & 12\,606 & 51\,603 \\
     & Samo instancjonowanie & 148 & 304 & 417 & 507 \\
    \midrule
    \multirow{3}{*}{\kol{Instanced}} & SRP Batcher + Resident Drawer & 216,5 & 432 & 213 & 267,5 \\
     & Sam SRP Batcher & 217 & 432 & 213 & 267,5 \\
     & Samo instancjonowanie & 217 & 432 & 213 & 267,5 \\
    \midrule
    \multirow{3}{*}{\kol{Combined}} & SRP Batcher + Resident Drawer & 39 & 122 & 225 & 613 \\
     & Sam SRP Batcher & 39 & 122 & 225 & 613 \\
     & Samo instancjonowanie & 39 & 122 & 225 & 613 \\
    \midrule
    \multirow{3}{*}{\kol{CombinedAtlas}} & SRP Batcher + Resident Drawer & 19 & 27 & 35 & 67 \\
     & Sam SRP Batcher & 19 & 27 & 35 & 67 \\
     & Samo instancjonowanie & 19 & 27 & 35 & 67 \\
    \bottomrule
  \end{tabular}
\end{table}

\CoTuMaByc{Z \kol{E7_*_warianty.csv} albo mediana z głównego pliku: \kol{narysowane_obiekty_mediana}.}

\begin{table}
  \centering
  \footnotesize
  \caption{E7: mediana narysowanych obiektów na wywołanie rysowania}
  \label{tab:e7-obiekty}
  \begin{tabular}{@{}llrrrr@{}}
    \toprule
    Rysowanie & Łączenie rysowania & 10$\times$10$\times$8 & 20$\times$20$\times$13 & 40$\times$40$\times$13 & 80$\times$80$\times$13 \\
    \midrule
    \multirow{3}{*}{\kol{GameObjects}} & SRP Batcher + Resident Drawer & 5,8 & 15,9 & 60,4 & 244,6 \\
     & Sam SRP Batcher & 1,0 & 1,0 & 1,0 & 1,0 \\
     & Samo instancjonowanie & 4,7 & 9,5 & 30,3 & 101,8 \\
    \midrule
    \multirow{3}{*}{\kol{Instanced}} & SRP Batcher + Resident Drawer & 6,1 & 16,7 & 59,8 & 194,0 \\
     & Sam SRP Batcher & 6,1 & 16,7 & 59,8 & 194,0 \\
     & Samo instancjonowanie & 6,1 & 16,7 & 59,8 & 194,0 \\
    \midrule
    \multirow{3}{*}{\kol{Combined}} & SRP Batcher + Resident Drawer & 1,0 & 1,0 & 1,0 & 1,0 \\
     & Sam SRP Batcher & 1,0 & 1,0 & 1,0 & 1,0 \\
     & Samo instancjonowanie & 1,0 & 1,0 & 1,0 & 1,0 \\
    \midrule
    \multirow{3}{*}{\kol{CombinedAtlas}} & SRP Batcher + Resident Drawer & 1,0 & 1,0 & 1,0 & 1,0 \\
     & Sam SRP Batcher & 1,0 & 1,0 & 1,0 & 1,0 \\
     & Samo instancjonowanie & 1,0 & 1,0 & 1,0 & 1,0 \\
    \bottomrule
  \end{tabular}
\end{table}

\CoTuMaByc{Z \kol{E7_*_warianty.csv} albo mediana z głównego pliku: \kol{setpass}.}

\begin{table}
  \centering
  \footnotesize
  \caption{E7: mediana zmian stanu SetPass na klatkę}
  \label{tab:e7-setpass}
  \begin{tabular}{@{}llrrrr@{}}
    \toprule
    Rysowanie & Łączenie rysowania & 10$\times$10$\times$8 & 20$\times$20$\times$13 & 40$\times$40$\times$13 & 80$\times$80$\times$13 \\
    \midrule
    \multirow{3}{*}{\kol{GameObjects}} & SRP Batcher + Resident Drawer & 9 & 9 & 9 & 9 \\
     & Sam SRP Batcher & 9 & 15 & 37 & 37 \\
     & Samo instancjonowanie & 45 & 61 & 83 & 83 \\
    \midrule
    \multirow{3}{*}{\kol{Instanced}} & SRP Batcher + Resident Drawer & 31 & 34 & 31 & 31 \\
     & Sam SRP Batcher & 31 & 34 & 31 & 31 \\
     & Samo instancjonowanie & 31 & 34 & 31 & 31 \\
    \midrule
    \multirow{3}{*}{\kol{Combined}} & SRP Batcher + Resident Drawer & 9 & 11 & 9 & 11 \\
     & Sam SRP Batcher & 9 & 11 & 9 & 9 \\
     & Samo instancjonowanie & 29 & 55 & 55 & 55 \\
    \midrule
    \multirow{3}{*}{\kol{CombinedAtlas}} & SRP Batcher + Resident Drawer & 9 & 11 & 9 & 9 \\
     & Sam SRP Batcher & 9 & 11 & 9 & 9 \\
     & Samo instancjonowanie & 9 & 11 & 9 & 9 \\
    \bottomrule
  \end{tabular}
\end{table}

\CoTuMaByc{Z \kol{E7_*_warianty.csv} albo mediana z głównego pliku: \kol{czas_klatki_ms_mediana}.}

\begin{table}
  \centering
  \footnotesize
  \caption{E7: mediana czasu klatki [ms]}
  \label{tab:e7-klatka}
  \begin{tabular}{@{}llrrrr@{}}
    \toprule
    Rysowanie & Łączenie rysowania & 10$\times$10$\times$8 & 20$\times$20$\times$13 & 40$\times$40$\times$13 & 80$\times$80$\times$13 \\
    \midrule
    \multirow{3}{*}{\kol{GameObjects}} & SRP Batcher + Resident Drawer & 0,513 & 0,478 & 0,552 & 0,691 \\
     & Sam SRP Batcher & 0,487 & 0,717 & 1,707 & 6,039 \\
     & Samo instancjonowanie & 0,509 & 0,747 & 1,446 & 4,252 \\
    \midrule
    \multirow{3}{*}{\kol{Instanced}} & SRP Batcher + Resident Drawer & 0,558 & 0,630 & 0,683 & 1,126 \\
     & Sam SRP Batcher & 0,520 & 0,625 & 0,624 & 1,046 \\
     & Samo instancjonowanie & 0,513 & 0,638 & 0,651 & 1,091 \\
    \midrule
    \multirow{3}{*}{\kol{Combined}} & SRP Batcher + Resident Drawer & 0,426 & 0,505 & 0,468 & 0,560 \\
     & Sam SRP Batcher & 0,370 & 0,411 & 0,457 & 0,508 \\
     & Samo instancjonowanie & 0,393 & 0,484 & 0,535 & 0,690 \\
    \midrule
    \multirow{3}{*}{\kol{CombinedAtlas}} & SRP Batcher + Resident Drawer & 0,432 & 0,444 & 0,453 & 0,514 \\
     & Sam SRP Batcher & 0,365 & 0,405 & 0,475 & 0,458 \\
     & Samo instancjonowanie & 0,364 & 0,415 & 0,469 & 0,461 \\
    \bottomrule
  \end{tabular}
\end{table}

\CoTuMaByc{Z \kol{E7_*_warianty.csv} albo mediana z głównego pliku: \kol{najdluzsza_klatka_ms_mediana}.}

\begin{table}
  \centering
  \footnotesize
  \caption{E7: mediana najdłuższej klatki [ms]}
  \label{tab:e7-najdluzsza}
  \begin{tabular}{@{}llrrrr@{}}
    \toprule
    Rysowanie & Łączenie rysowania & 10$\times$10$\times$8 & 20$\times$20$\times$13 & 40$\times$40$\times$13 & 80$\times$80$\times$13 \\
    \midrule
    \multirow{3}{*}{\kol{GameObjects}} & SRP Batcher + Resident Drawer & 1,38 & 1,20 & 1,35 & 1,80 \\
     & Sam SRP Batcher & 1,20 & 1,04 & 2,17 & 8,08 \\
     & Samo instancjonowanie & 1,20 & 1,14 & 1,90 & 6,26 \\
    \midrule
    \multirow{3}{*}{\kol{Instanced}} & SRP Batcher + Resident Drawer & 1,27 & 1,23 & 1,32 & 2,03 \\
     & Sam SRP Batcher & 1,21 & 1,19 & 1,32 & 1,93 \\
     & Samo instancjonowanie & 1,21 & 1,20 & 1,31 & 2,12 \\
    \midrule
    \multirow{3}{*}{\kol{Combined}} & SRP Batcher + Resident Drawer & 1,13 & 1,09 & 1,23 & 1,16 \\
     & Sam SRP Batcher & 1,15 & 1,21 & 1,23 & 1,18 \\
     & Samo instancjonowanie & 1,17 & 1,23 & 1,25 & 1,20 \\
    \midrule
    \multirow{3}{*}{\kol{CombinedAtlas}} & SRP Batcher + Resident Drawer & 1,14 & 1,18 & 1,21 & 1,27 \\
     & Sam SRP Batcher & 1,14 & 1,16 & 1,38 & 1,28 \\
     & Samo instancjonowanie & 1,17 & 1,15 & 1,27 & 1,31 \\
    \bottomrule
  \end{tabular}
\end{table}

\CoTuMaByc{Z \kol{E7_*_warianty.csv}: \kol{czas_stawiania_ms_mediana}. Podaj, dla którego ustawienia łączenia rysowania.}

\begin{table}
  \centering
  \footnotesize
  \caption[E7: mediana czasu stawiania mapy]{E7: mediana czasu stawiania mapy [ms], sam SRP Batcher}
  \label{tab:e7-stawianie}
  \begin{tabular}{@{}lrrrr@{}}
    \toprule
    Rysowanie & 10$\times$10$\times$8 & 20$\times$20$\times$13 & 40$\times$40$\times$13 & 80$\times$80$\times$13 \\
    \midrule
    \kol{GameObjects} & 0,48 & 2,08 & 10,98 & 62,10 \\
    \kol{Instanced} & 0,15 & 0,33 & 0,85 & 3,25 \\
    \kol{Combined} & 0,47 & 1,15 & 3,53 & 11,14 \\
    \kol{CombinedAtlas} & 0,52 & 1,07 & 3,28 & 10,38 \\
    \bottomrule
  \end{tabular}
\end{table}

\CoTuMaByc{Z \kol{E7_*_warianty.csv}: \kol{pamiec_unity_mb_mediana}. Podaj, dla którego ustawienia łączenia rysowania.}

\begin{table}
  \centering
  \footnotesize
  \caption[E7: mediana pamięci silnika]{E7: mediana pamięci silnika [MB], sam SRP Batcher}
  \label{tab:e7-pamiec}
  \begin{tabular}{@{}lrrrr@{}}
    \toprule
    Rysowanie & 10$\times$10$\times$8 & 20$\times$20$\times$13 & 40$\times$40$\times$13 & 80$\times$80$\times$13 \\
    \midrule
    \kol{GameObjects} & 303,8 & 313,1 & 358,7 & 397,9 \\
    \kol{Instanced} & 302,6 & 307,7 & 353,3 & 368,7 \\
    \kol{Combined} & 297,0 & 308,8 & 370,1 & 425,7 \\
    \kol{CombinedAtlas} & 297,0 & 308,7 & 370,1 & 425,9 \\
    \bottomrule
  \end{tabular}
\end{table}

\CoTuMaByc{Z głównego pliku E7: \kol{pamiec_siatek_mb}, \kol{bok_fragmentu}.}

\begin{table}
  \centering
  \footnotesize
  \caption[E7: rozmiar scalonych siatek]{E7: mediana rozmiaru scalonych siatek [MB], fragmenty o boku 16 komórek}
  \label{tab:e7-siatki}
  \begin{tabular}{@{}lrrrr@{}}
    \toprule
    Rysowanie & 10$\times$10$\times$8 & 20$\times$20$\times$13 & 40$\times$40$\times$13 & 80$\times$80$\times$13 \\
    \midrule
    \kol{Combined} & 0,55 & 2,31 & 10,57 & 44,96 \\
    \kol{CombinedAtlas} & 0,55 & 2,31 & 10,57 & 44,96 \\
    \bottomrule
  \end{tabular}
\end{table}

\CoTuMaByc{Z \kol{E8_*_warianty.csv}: \kol{gestosc_regul_proc}, \kol{gestosc_wazona_proc}, \kol{kappa}, \kol{udane_proby_proc}, \kol{zabudowa_proc}.}

\begin{table}
  \centering
  \footnotesize
  \caption{E8: restrykcyjność reguł i powodzenie generacji}
  \label{tab:e8}
  \begin{tabularx}{\linewidth}{@{}lCCCCCC@{}}
    \toprule
    Zestaw lub wariant & Stany & Gęstość par [\%] & Gęstość ważona [\%] & $\kappa$ & Udane próby [\%] & Zabudowa [\%] \\
    \midrule
    \kol{Prosty} & 26 & 7,79 & 21,62 & 2,350 & 46,1 & 29,1 \\
    \kol{Sredni} & 79 & 10,24 & 18,99 & 1,565 & 79,4 & 35,8 \\
    \kol{PelnyBezDekoracji} & 146 & 8,45 & 12,10 & 1,487 & 62,5 & 40,5 \\
    \kol{Pelny} & 170 & 8,05 & 11,67 & 1,471 & 71,9 & 39,4 \\
    \midrule
    \kol{PelnyBezWolnegoPola} & 170 & 13,37 & 15,65 & 1,175 & 9,5 & 56,6 \\
    \kol{PelnyBezDrzwiNadSchodami} & 170 & 8,17 & 11,79 & 1,463 & 88,5 & 40,4 \\
    \kol{PelnyTrawa040} & 170 & 8,05 & 11,06 & 1,471 & 46,9 & 50,6 \\
    \kol{PelnyTrawa025} & 170 & 8,05 & 10,93 & 1,471 & 32,3 & 55,1 \\
    \bottomrule
  \end{tabularx}
\end{table}

\CoTuMaByc{Z \kol{*_korelacje.csv}: \kol{miara}, \kol{warianty}, \kol{liczba_wariantow}, \kol{spearman_z_udanymi_probami}. Mapa zdegenerowana w kodzie: zabudowa parteru poniżej 1\% albo powyżej 90\%. Korelacja rang: \cite{Spearman1904}.}

\begin{table}
  \centering
  \footnotesize
  \caption{Korelacja Spearmana restrykcyjności z odsetkiem udanych prób}
  \label{tab:korelacje}
  \begin{tabular}{@{}llrrrr@{}}
    \toprule
    Miara & Warianty & E8: liczba & E8: $\rho$ & E9: liczba & E9: $\rho$ \\
    \midrule
    \multirow{2}{*}{Gęstość par} & wszystkie & 8 & 0,098 & 81 & 0,388 \\
     & bez map zdegenerowanych & 8 & 0,098 & 75 & 0,531 \\
    \midrule
    \multirow{2}{*}{Gęstość ważona} & wszystkie & 8 & 0,048 & 81 & 0,430 \\
     & bez map zdegenerowanych & 8 & 0,048 & 75 & 0,569 \\
    \midrule
    \multirow{2}{*}{$\kappa$} & wszystkie & 8 & 0,146 & 81 & $-$0,388 \\
     & bez map zdegenerowanych & 8 & 0,146 & 75 & $-$0,531 \\
    \bottomrule
  \end{tabular}
\end{table}

\CoTuMaByc{Z \kol{E9_*_warianty.csv}: \kol{zaostrzenie_proc}, \kol{losowanie_zaostrzenia}, \kol{udane_proby_proc}, \kol{zabudowa_proc}.}

\begin{table}
  \centering
  \footnotesize
  \caption{E9: wyniki według poziomu zaostrzenia reguł}
  \label{tab:e9}
  \begin{tabularx}{\linewidth}{@{}CCCCCCC@{}}
    \toprule
    Zaostrzenie [\%] & Losowania & Udane próby: mediana [\%] & Udane próby: minimum [\%] & Udane próby: maksimum [\%] & Zabudowa: mediana [\%] & Warianty z mapami zdegenerowanymi \\
    \midrule
    0 & 1 & 71,4 & 71,4 & 71,4 & 39,6 & 0 \\
    0,2 & 10 & 64,5 & 58,8 & 74,1 & 39,9 & 0 \\
    0,5 & 10 & 31,7 & 0,8 & 66,7 & 40,1 & 0 \\
    1 & 10 & 30,0 & 0,4 & 66,7 & 38,5 & 1 \\
    2 & 10 & 7,9 & 0,0 & 76,9 & 35,6 & 1 \\
    3 & 10 & 6,0 & 0,0 & 60,6 & 36,9 & 1 \\
    5 & 10 & 17,7 & 0,0 & 87,0 & 26,7 & 0 \\
    8 & 10 & 11,3 & 0,0 & 100,0 & 3,7 & 1 \\
    12 & 10 & 1,0 & 0,0 & 100,0 & 0,0 & 2 \\
    \bottomrule
  \end{tabularx}
\end{table}

\CoTuMaByc{Z głównego pliku E1: \kol{decyzje}, \kol{kroki_propagacji}, \kol{usuniete_propagacja}, \kol{usuniete_decyzje}. Udział i usunięcia na decyzję: \cref{eq:propagacja}.}

\begin{table}
  \centering
  \footnotesize
  \caption[Propagacja według rozwiązywacza (E1)]{Propagacja według rozwiązywacza (E1), mediany na mapę}
  \label{tab:propagacja-rozwiazywacze}
  \begin{tabularx}{\linewidth}{@{}lCCCCCC@{}}
    \toprule
    Rozwiązywacz & Decyzje & Kroki propagacji & Usunięte przez propagację & Usunięte przez decyzje & Udział propagacji [\%] & Usunięcia na decyzję \\
    \midrule
    \kol{Naive} & 702 & 273\,449 & 684\,265 & 2564 & 99,66 & 1092,3 \\
    \kol{SupportCount} & 702 & 4\,030\,944 & 684\,265 & 2564 & 99,66 & 1092,3 \\
    \kol{Bitset} & 702 & 273\,449 & 684\,265 & 2564 & 99,66 & 1092,3 \\
    \kol{Burst} & 702 & 273\,449 & 684\,265 & 2564 & 99,66 & 1092,3 \\
    \bottomrule
  \end{tabularx}
\end{table}

\CoTuMaByc{Z głównego pliku E4; podaj, dla którego rozwiązywacza.}

\begin{table}
  \centering
  \footnotesize
  \caption[Propagacja według rozmiaru siatki (E4)]{Propagacja według rozmiaru siatki (E4), rozwiązywacz Bitset, mediany na mapę. Liczniki obejmują wszystkie próby, razem z restartami}
  \label{tab:propagacja-rozmiary}
  \begin{tabularx}{\linewidth}{@{}lCCCCCC@{}}
    \toprule
    Siatka & Decyzje & Kroki propagacji & Usunięte przez propagację & Usunięte przez decyzje & Udział propagacji [\%] & Usunięcia na decyzję \\
    \midrule
    10$\times$10$\times$8 & 160 & 35\,891 & 84\,153 & 601 & 99,33 & 566,5 \\
    20$\times$20$\times$13 & 702 & 273\,449 & 684\,265 & 2564 & 99,66 & 1092,3 \\
    40$\times$40$\times$13 & 5504 & 3\,436\,857 & 10\,646\,474 & 25\,646 & 99,76 & 2011,5 \\
    80$\times$80$\times$13 & 991\,951 & 1\,233\,864\,193 & 5\,603\,406\,919 & 5\,882\,402 & 99,89 & 5490,4 \\
    \bottomrule
  \end{tabularx}
\end{table}

\CoTuMaByc{Kolumny głównych plików: \kol{udane}, \kol{naruszenia}.}

\begin{table}
  \centering
  \footnotesize
  \caption{Poprawność map we wszystkich eksperymentach}
  \label{tab:naruszenia}
  \begin{tabular}{@{}lrrr@{}}
    \toprule
    Eksperyment & Uruchomienia & Udane mapy & Naruszenia zasad \\
    \midrule
    E1 & 400 & 400 & 0 \\
    E2 & 500 & 500 & 0 \\
    E3 & 1200 & 1150 & 0 \\
    E4 & 1600 & 1350 & 0 \\
    E5 & 400 & 400 & 0 \\
    E6 & 500 & 500 & 0 \\
    E7 & 1440 & 1440 & 0 \\
    E8 & 800 & 800 & 0 \\
    E9 & 1620 & 1100 & 0 \\
    \bottomrule
  \end{tabular}
\end{table}


\section{Materiał na załączniki}
%%% KLOCKI: MATERIAŁ NA ZAŁĄCZNIKI
%%% Dane z kodu: CsvWriter.cs, BenchmarkSession.cs, BenchmarkRunner.cs. Załącznik zakładasz przez \chapter po \appendix w main.tex.

\CoTuMaByc{Kolejność jak w pliku. Kolumny związane z rysowaniem są puste bez okna.}

\begin{longtable}{@{}l>{\raggedright\arraybackslash}p{8.9cm}l@{}}
  \caption{Kolumny pliku wyników}\label{tab:kolumny-csv}\\
  \toprule
    Kolumna & Znaczenie & Jednostka \\
    \midrule
  \endfirsthead
  \caption[]{Kolumny pliku wyników (ciąg dalszy)}\\
  \toprule
    Kolumna & Znaczenie & Jednostka \\
    \midrule
  \endhead
  \bottomrule
  \endlastfoot
    \kol{data} & data i godzina zapisu wiersza & -- \\
    \kol{komputer} & nazwa komputera & -- \\
    \kol{procesor} & model procesora & -- \\
    \kol{ram_gb} & pamięć operacyjna & GB \\
    \kol{wersja_unity} & wersja Unity & -- \\
    \kol{wersja_programu} & wersja aplikacji & -- \\
    \kol{eksperyment} & nazwa eksperymentu & -- \\
    \kol{rozwiazywacz} & rozwiązywacz & -- \\
    \kol{kolejnosc} & heurystyka wyboru komórki & -- \\
    \kol{sprzecznosc} & reakcja na sprzeczność & -- \\
    \kol{zestaw} & zestaw kafelków albo wariant reguł & -- \\
    \kol{x} & rozmiar siatki: E & komórki \\
    \kol{y} & rozmiar siatki: N & komórki \\
    \kol{z} & rozmiar siatki: piętra & komórki \\
    \kol{ziarno} & ziarno generatora liczb losowych & -- \\
    \kol{powtorzenie} & numer powtórzenia ziarna & -- \\
    \kol{udane} & czy mapa powstała: 1 albo 0 & -- \\
    \kol{proby} & liczba prób & -- \\
    \kol{czas_ms} & czas generacji & ms \\
    \kol{czas_wybor_ms} & czas wyboru komórek & ms \\
    \kol{czas_propagacja_ms} & czas propagacji & ms \\
    \kol{czas_cofanie_ms} & czas cofania & ms \\
    \kol{decyzje} & liczba decyzji & -- \\
    \kol{kroki_propagacji} & liczba kroków propagacji & -- \\
    \kol{usuniete_mozliwosci} & wszystkie usunięte możliwości & -- \\
    \kol{sprzecznosci} & liczba sprzeczności & -- \\
    \kol{cofniecia} & liczba cofnięć & -- \\
    \kol{restarty} & liczba startów od nowa & -- \\
    \kol{pamiec_szczytowa_mb} & największy rozmiar struktur rozwiązywacza i selektora, szacowany z rozmiarów tablic & MB \\
    \kol{alokacje_gc_kb} & pamięć zaalokowana na stercie przez rozwiązywacz; w aplikacji przyrost sterty & KB \\
    \kol{naruszenia} & naruszenia zasad w gotowej mapie & -- \\
    \kol{wieze} & liczba wież & -- \\
    \kol{donzony} & liczba donżonów & -- \\
    \kol{mosty} & liczba mostów & -- \\
    \kol{schody} & liczba schodów & -- \\
    \kol{wys_wiez_sr} & średnia wysokość wież & piętra \\
    \kol{wys_murow_sr} & średnia wysokość murów & piętra \\
    \kol{wys_donzonow_sr} & średnia wysokość donżonów & piętra \\
    \kol{trawa_dostepna_proc} & odsetek terenu na parterze połączonego z krawędzią mapy & \% \\
    \kol{dziedzince_zamkniete} & obszary terenu na parterze bez połączenia z krawędzią mapy & -- \\
    \kol{czas_stawiania_ms} & czas stawiania mapy w oknie & ms \\
    \kol{srodowisko} & rodzaj uruchomienia: Edytor, Mono albo IL2CPP, z dopiskiem przy wersji deweloperskiej & -- \\
    \kol{roznica_mapy_proc} & odsetek komórek innych niż w mapie poprzedniego ziarna: wszystkie komórki & \% \\
    \kol{roznica_parter_proc} & to samo: parter & \% \\
    \kol{roznica_zabudowa_proc} & to samo: komórki zabudowy & \% \\
    \kol{rysowanie} & sposób rysowania & -- \\
    \kol{wywolania_rysowania} & wywołania rysowania na klatkę, razem z cieniami i oknem ustawień & -- \\
    \kol{narysowane_obiekty} & obiekty narysowane przez te wywołania & -- \\
    \kol{setpass} & zmiany stanu SetPass na klatkę & -- \\
    \kol{trojkaty} & trójkąty na klatkę & -- \\
    \kol{wierzcholki} & wierzchołki na klatkę & -- \\
    \kol{czas_klatki_ms} & średni czas między klatkami & ms \\
    \kol{czas_klatki_cpu_ms} & czas pracy procesora nad klatką & ms \\
    \kol{czas_klatki_gpu_ms} & czas pracy karty graficznej nad klatką & ms \\
    \kol{usuniete_propagacja} & możliwości usunięte przez propagację & -- \\
    \kol{usuniete_decyzje} & możliwości usunięte przez decyzje & -- \\
    \kol{pamiec_rzeczywista_mb} & przyrost sterty po generacji i odśmieceniu, gdy rozwiązywacz jeszcze żyje & MB \\
    \kol{zabudowa_proc} & odsetek parteru zajęty przez budynki & \% \\
    \kol{gestosc_regul_proc} & gęstość par & \% \\
    \kol{gestosc_wazona_proc} & gęstość ważona & \% \\
    \kol{kappa} & współczynnik $\kappa$ & -- \\
    \kol{zaostrzenie_proc} & odsetek usuniętych dozwolonych par & \% \\
    \kol{losowanie_zaostrzenia} & numer losowania usuwanych par & -- \\
    \kol{laczenie_rysowania} & ustawienie łączenia rysowania w URP & -- \\
    \kol{bok_fragmentu} & bok fragmentu scalonej mapy; puste w innych sposobach rysowania & komórki \\
    \kol{pamiec_siatek_mb} & rozmiar scalonych siatek & MB \\
    \kol{pamiec_unity_mb} & pamięć przydzielona przez silnik & MB \\
    \kol{pamiec_gc_mb} & sterta zarządzana & MB \\
    \kol{pamiec_procesu_mb} & pamięć procesu w systemie & MB \\
    \kol{najdluzsza_klatka_ms} & najdłuższa klatka z pomiaru & ms \\
    \kol{pierwsza_klatka_ms} & od postawienia mapy do początku następnej klatki & ms \\
\end{longtable}

\begin{table}
  \centering
  \footnotesize
  \caption{Pliki z podsumowaniem serii}
  \label{tab:pliki-podsumowania}
  \begin{tabularx}{\linewidth}{@{}lL@{}}
    \toprule
    Plik & Kolumny \\
    \midrule
    \kol{*_warianty.csv} & \kol{eksperyment}, \kol{rozwiazywacz}, \kol{kolejnosc}, \kol{sprzecznosc}, \kol{zestaw}, \kol{x}, \kol{y}, \kol{z}, \kol{rysowanie}, \kol{laczenie_rysowania}, \kol{zaostrzenie_proc}, \kol{losowanie_zaostrzenia}, \kol{mapy}, \kol{udane_mapy}, \kol{proby}, \kol{sprzecznosci}, \kol{udane_proby_proc}, \kol{zabudowa_proc}, \kol{gestosc_regul_proc}, \kol{gestosc_wazona_proc}, \kol{kappa}, \kol{czas_ms_mediana}, \kol{czas_ms_q1}, \kol{czas_ms_q3}, \kol{czas_stawiania_ms_mediana}, \kol{czas_klatki_ms_mediana}, \kol{najdluzsza_klatka_ms_mediana}, \kol{wywolania_rysowania_mediana}, \kol{narysowane_obiekty_mediana}, \kol{pamiec_rzeczywista_mb_mediana}, \kol{pamiec_unity_mb_mediana} \\
    \kol{*_korelacje.csv} & \kol{miara}, \kol{warianty}, \kol{liczba_wariantow}, \kol{spearman_z_udanymi_probami} \\
    \bottomrule
  \end{tabularx}
\end{table}

\begin{table}
  \centering
  \footnotesize
  \caption{Argumenty uruchomienia aplikacji}
  \label{tab:argumenty}
  \begin{tabularx}{\linewidth}{@{}lL@{}}
    \toprule
    Argument & Znaczenie \\
    \midrule
    \kol{-benchmark} & badanie startuje samo, a aplikacja zamyka się po nim \\
    \kol{-config E1} & eksperyment do uruchomienia; \kol{-config wszystkie} uruchamia E1--E9 po kolei \\
    \kol{-out plik.csv} & plik wyników; domyślnie folder \kol{Wyniki} obok aplikacji \\
    \kol{-seeds N}, \kol{-firstSeed N} & liczba ziaren i pierwsze ziarno \\
    \kol{-repetitions N}, \kol{-warmup N} & powtórzenia ziarna i uruchomienia na rozgrzewkę \\
    \kol{-showEvery N} & w oknie pokazuje co N-tą mapę \\
    \kol{-display}, \kol{-batching} & sposób rysowania i łączenie rysowania \\
    \kol{-batchmode -nographics} & bez okna i bez pomiaru rysowania; E7 działa tylko w oknie \\
    \kol{-screen-fullscreen 0} & okno zamiast pełnego ekranu; rozmiar: \kol{-screen-width}, \kol{-screen-height} \\
    \bottomrule
  \end{tabularx}
\end{table}

\CoTuMaByc{Pierwsze polecenie to zaplanowana seria pomiarów, drugie to przykład trybu bez okna. Po pomiarach popraw je na takie, jakich naprawdę użyto.}

\begin{lstlisting}[numbers=none, caption={Przykładowe uruchomienie badań}, label={lst:uruchomienie}]
./Build/Zamek.exe -screen-fullscreen 0 -screen-width 1600 -screen-height 900 -benchmark -config wszystkie -seeds 30
./Build/Zamek.exe -batchmode -nographics -benchmark -config E1 -out wyniki_E1.csv
\end{lstlisting}

\CoTuMaByc{Adres repozytorium: \url{https://github.com/zaaabinski/Inzynierka}.}

\begin{table}
  \centering
  \footnotesize
  \caption{Zawartość repozytorium}
  \label{tab:repozytorium}
  \begin{tabularx}{\linewidth}{@{}lL@{}}
    \toprule
    Folder & Zawartość \\
    \midrule
    \kol{Assets/Scripts/Core} & reguły, rozwiązywacze, analiza map \\
    \kol{Assets/Scripts/Runtime} & zasoby, stawianie mapy, badania, panel aplikacji \\
    \kol{Assets/Scripts/Editor} & okno generatora, import reguł i modeli, budowanie aplikacji \\
    \kol{Assets/Scripts/Tests} & testy automatyczne \\
    \kol{Assets/Data/Rules} & plik reguł \\
    \kol{Assets/Data/Tiles}, \kol{Assets/Data/Tilesets} & kafelki i zestawy \\
    \kol{Assets/Data/Experiments} & ustawienia eksperymentów E1--E9 \\
    \kol{Assets/Materials}, \kol{Assets/Prefabs}, \kol{Assets/Scenes} & materiały, prefaby, sceny \\
    \kol{Analiza} & skrypt analizy wyników \\
    \kol{Wyniki} & pliki wyników \\
    \kol{Praca} & tekst pracy w LaTeX-u \\
    \bottomrule
  \end{tabularx}
\end{table}


\endgroup

%% Propozycje literatury widać w szkicu; po \szkicfalse w spisie zostaje tylko to, co zacytowane.
\nocite{*}

\fi

%% Wskazówka pochodzi z wymagań Wydziału dla pracy inżynierskiej.

\chapter{Zakończenie}
\label{chap:zakonczenie}

\CoTuMaByc{Wymagania Wydziału: najważniejsze rezultaty i wnioski; odniesienie do celów ze wstępu (czy zrealizowano w pełni, a jeśli nie, to dlaczego); ograniczenia oraz mocne i słabe strony rozwiązania; możliwości dalszego rozwinięcia; obszar praktycznego wykorzystania wyników.}


%% BIBLIOGRAFIA
\printbibliography[heading=bibintoc, title={Spis literatury}]

%% ZAŁĄCZNIKI
\appendix
%% % !TEX root = ../main.tex
%% Załącznik A: pliki wyników i uruchomienie aplikacji.

\chapter{Pliki wyników i uruchomienie badań}
\label{app:pliki}

\CoTuMaByc{Klocki z \kol{klocki/zalaczniki.tex}: \kol{tab:kolumny-csv}, \kol{tab:pliki-podsumowania}, \kol{tab:argumenty}, \kol{lst:uruchomienie}, \kol{tab:repozytorium}.}

%% \include{chapters/zalacznik_b}

%% SPISY
\clearpage
\phantomsection
\addcontentsline{toc}{chapter}{\listfigurename}
\listoffigures
\clearpage
\phantomsection
\addcontentsline{toc}{chapter}{\listtablename}
\listoftables
\clearpage
\phantomsection
\addcontentsline{toc}{chapter}{\lstlistlistingname}
\lstlistoflistings

\end{document}
